\section{综合案例、实验与源码对照}

这一章专门补齐三个容易被刷题书忽略的环节：第一，经典题不能直接给最终答案，而要从暴力方法、瓶颈定位、结构观察和正确性证明一步步推出来；第二，C++ 容器不是语法糖，容器的内存模型、迭代器失效、比较器和分配成本会改变工程质量；第三，Linux 源码里的数据结构不是孤立技巧，而是访问模式、生命周期、并发和内存分配共同作用的结果。读本章时，不要把它当题号清单，而要把它当作一组可以反复做的实验。

\subsection{经典案例推导}

\noindent\textbf{经典题完整推导。}

\noindent\textbf{LeetCode 1 两数之和。}

这道题看似简单，却是“从暴力到结构”的最小样本。题面要求找两个下标，值相加等于目标。最直接的暴力做法是两层循环枚举 \code{i} 和 \code{j}，检查 \code{nums[i] + nums[j] == target}。这个版本一定正确，因为它覆盖了所有下标对；它的问题是同一个历史值会被未来很多位置反复扫描。把瓶颈说清楚后，优化就自然出现：当扫描到当前值 \code{x} 时，需要的另一个数只可能是 \code{target - x}，所以不必再遍历所有历史下标，只要查历史表即可。

这里的哈希表不是“模板”，而是把暴力中的历史扫描换成等值查询。表的 key 是已经出现过的数，value 是它的下标。顺序必须先查再插入当前数，因为题目不允许同一个下标使用两次。如果先插入再查，输入 \code{[3]}、目标 \code{6} 这种边界就会把当前元素和自己配对。重复数字也要靠这个顺序处理：输入 \code{[3,3]}、目标 \code{6} 时，扫描第二个 \code{3} 才能查到第一个 \code{3}。

正确性可以对照暴力证明。任意合法答案有两个下标 \code{i < j}。算法扫描到 \code{j} 时，\code{i} 的值已经插入历史表；由于 \code{nums[i] + nums[j] == target}，查 \code{target - nums[j]} 一定能找到 \code{i}。反过来，算法找到的历史下标一定早于当前下标，且值满足目标和，所以返回结果合法。复杂度从 \complexity{O(n^2)} 降到平均 \complexity{O(n)}，代价是 \complexity{O(n)} 额外空间。实验时要覆盖没有答案、重复值、负数、目标为零、补数等于当前值这些输入。

\noindent\textbf{LeetCode 3 无重复字符的最长子串。}

暴力版本可以枚举每个左端和右端，再检查子串内是否有重复字符。如果检查也用集合扫描，复杂度会很高。优化的第一步不是立刻说“滑动窗口”，而是观察连续子串的变化：右端向右扩展时，只新增一个字符；如果这个字符在当前窗口里重复，左端只需要移动到它上一次出现位置之后。左端永远不会向左回退，因为已经被排除的左端只会让窗口更长、更容易包含重复字符。

高质量实现通常保存字符上一次出现的位置。扫描右端 \code{right} 时，如果字符 \code{s[right]} 的上次位置在当前窗口内部，就把 \code{left} 更新为 \code{last[s[right]] + 1}；如果上次位置在窗口左侧之外，就不能把 \code{left} 拉回去。这里最常见的错误是直接写 \code{left = last[c] + 1}，导致输入 \code{abba} 时，处理最后一个 \code{a} 会把左端从 2 错误拉回 1。

这题的正确性来自窗口不变量：任意时刻，\code{[left,right]} 内没有重复字符，且 \code{left} 是在当前 \code{right} 下能保持无重复的最左合法边界。每次右端加入一个字符，只可能因为这个新字符破坏合法性；把左端移过它上次出现位置后，窗口重新合法。答案在每次窗口合法后更新。实验应打印 \code{left}、\code{right}、当前字符、上次位置和答案，特别观察 \code{abba}、空串、全相同字符、全部不同字符以及扩展字符集的情况。

\noindent\textbf{LeetCode 15 三数之和。}

暴力方法三层循环枚举三个下标，能够覆盖全部答案，但复杂度是 \complexity{O(n^3)}，并且重复组合处理很麻烦。优化从排序开始。排序的意义不是让数组变好看，而是制造单调性：固定第一个数以后，剩下两个数在有序数组中寻找目标和。如果左指针和右指针的和太小，左指针右移才能增大；如果和太大，右指针左移才能减小。这样一层固定加一层双指针就能覆盖所有二元组合。

去重是这题的核心边界。固定第一个数时，如果它和前一个固定值相同，就跳过，否则会生成重复三元组。找到一个三元组后，左右指针也要分别跳过相同值，避免同一组值被多个下标组合重复加入。注意，去重不是随便用集合过滤结果；集合过滤可以让答案看起来对，但会掩盖推导过程，也可能带来额外排序和哈希成本。

正确性可以从固定值证明。排序后，对于每个固定位置 \code{i}，双指针会在区间 \code{(i,n)} 中按单调规则排除不可能的配对。若当前和小于目标，所有使用当前左指针和更小右指针的组合只会更小，所以左指针可以安全右移；若当前和大于目标，所有使用当前右指针和更大左指针的组合只会更大，所以右指针可以安全左移。实验要覆盖全零、大量重复、无解、正负混合、答案很多以及目标和可能溢出的四数扩展。

\noindent\textbf{LeetCode 42 接雨水。}

这题最容易被误讲成“背双指针”。先从暴力出发：对每个位置 \code{i}，它能接的水量由左侧最高柱和右侧最高柱的较小者决定，再减去当前高度。暴力每个位置都向两边扫描，复杂度 \complexity{O(n^2)}。第一层优化是预处理左右最高数组，令 \code{left_max[i]} 表示 \code{i} 左侧含自身最高，\code{right_max[i]} 表示右侧含自身最高。这样每个位置的贡献能常数时间计算。

双指针进一步把左右最高数组压缩成两个变量。左指针和右指针从两端向中间走，同时维护 \code{left_max} 和 \code{right_max}。如果 \code{left_max <= right_max}，左侧当前位置的水位已经由 \code{left_max} 决定，因为右侧至少存在一个不低于 \code{left_max} 的边界；此时可以结算左指针并右移。反之结算右指针。这个证明比代码更重要：结算的一侧必须已经知道较小边界，否则会提前计算。

单调栈解法则从另一角度结算凹槽。栈中保存递减高度下标，遇到更高的右边界时，弹出槽底，用当前柱作为右边界、新栈顶作为左边界计算水量。三种解法都正确，但状态不同：前后缀数组保存每个位置的边界，双指针保存当前两侧边界，单调栈保存尚未找到右边界的柱子。实验时分别打印三个版本的状态，对比单调递增、单调递减、平台高度、多个凹槽和长度不足三的情况。

\noindent\textbf{LeetCode 76 最小覆盖子串。}

暴力方法枚举所有子串，再检查它是否覆盖目标串字符需求。它的问题有两层：枚举窗口多，检查窗口也重复统计。优化要抓住覆盖关系的单调性。右端扩展时，窗口字符只会增加，覆盖条件更容易满足；一旦满足覆盖，继续扩展只会让窗口更长，不可能改善当前左端下的答案，所以应该尽量收缩左端，直到刚好失去覆盖。

实现中需要两个计数结构：目标需求 \code{need} 和窗口计数 \code{window}。为了避免每次比较全部字符，可以维护 \code{formed}，表示有多少种字符已经达到需求次数。注意这里是“达到需求次数的字符种类”，不是窗口里匹配的字符总数。目标串包含重复字符时尤其容易错，例如 \code{AABC} 需要两个 \code{A}，窗口里一个 \code{A} 不算满足。

正确性来自最短窗口的收缩。每当右端使窗口满足覆盖，内层循环会尽可能移动左端，并在每个合法窗口上更新答案。对于某个固定右端，循环结束前已经检查了所有能覆盖目标的左端中最短的那个；对于所有右端扫描完成后，全局最短自然被覆盖。实验要覆盖目标字符缺失、目标含重复、最短窗口在开头、最短窗口在结尾、大小写混合和源串短于目标串。C++ 中若字符集固定，可用数组；若字符集不确定，用哈希表但要注意 \code{operator[]} 的默认插入。

\noindent\textbf{LeetCode 84 柱状图中最大的矩形。}

暴力方法枚举每根柱子作为矩形高度，再向左右扩展直到遇到更矮柱子。这个做法清楚地暴露了瓶颈：每根柱子都在寻找左右第一个更矮的位置。单调栈优化就是一次扫描求出“某根柱子作为最矮高度时的最大宽度”。栈中保存高度递增的下标；当遇到一个更矮柱子时，栈顶柱子的右边界已经确定为当前位置，左边界是弹出后新的栈顶。

宽度计算是这题最容易错的地方。弹出下标 \code{mid} 后，如果栈为空，说明它左侧没有更矮柱子，宽度是 \code{right}；如果栈不空，左边界是 \code{stack.top()}，宽度是 \code{right - stack.top() - 1}。为了统一结算，可以在末尾补一个高度为零的哨兵，让所有柱子最终都被弹出。相等高度如何处理也要固定：常见写法遇到小于当前栈顶才弹，或者遇到小于等于时弹，只要宽度语义和去重一致即可。

正确性来自结算时机。某根柱子被弹出，说明当前柱子是它右侧第一个更矮位置；弹出后新的栈顶是它左侧第一个更矮位置。两个边界之间所有柱子高度都不低于它，所以以它为最矮柱的最大矩形已经确定，未来再向右也会被当前更矮柱阻断。实验要覆盖递增、递减、全相等、单柱、空数组、哨兵和宽度计算。这个题也是 LeetCode 85 最大矩形的基础：二维矩阵每一行都能转成一组柱状图高度。

\noindent\textbf{LeetCode 146 LRU 缓存。}

LRU 不是“哈希表加链表”六个字。题目要求 \code{get} 和 \code{put} 都是常数时间，并且容量满时淘汰最久未使用项。先看暴力：用数组或向量保存访问顺序，每次访问都线性查找并移动到头部；定位和移动都会慢。哈希表能按 key 常数定位，但无法维护新旧顺序；双向链表能常数删除已知节点并移动到头部，但无法按 key 查找。两个结构职责互补，组合起来才满足要求。

高质量实现要先定义节点身份。链表节点里保存 key 和 value；哈希表从 key 映射到链表迭代器或节点指针。为什么节点里还要保存 key？因为淘汰尾部节点时，需要从哈希表删除对应 key。如果节点只保存 value，尾部删除后无法常数时间更新哈希表。访问已有 key 时，更新 value 后要把节点移动到头部；插入新 key 时，容量超过上限就删除尾部。

正确性来自两个不变量：哈希表中的每个 key 都指向链表中的一个节点，链表从头到尾按最近使用到最久未使用排序。每次 \code{get} 或成功更新都会把节点移到头部，所以尾部始终是最久未使用项。实验要覆盖容量为一、重复 \code{put}、访问不存在 key、更新已有值、连续淘汰和插入后立即访问。对照 Linux 源码时，可以进一步理解侵入式链表：业务对象自己带链表节点，哈希索引只保存对象位置，链表不拥有对象。

\noindent\textbf{LeetCode 200 岛屿数量。}

这题的第一步是建模：网格中的陆地格子是节点，四方向相邻陆地之间有边，岛屿就是连通分量。暴力错误做法常常从每个陆地都重新搜索，或者不标记访问导致重复计数。正确方法是扫描每个格子，遇到一个未访问陆地，就启动 BFS 或 DFS，把整个连通块标记为已访问，并把岛屿数量加一。

BFS 和 DFS 在这题里都可以。BFS 用队列保存待扩展格子，入队时标记访问，避免同一格子被多个邻居重复加入；DFS 可以递归，也可以用显式栈。工程上大网格可能导致递归栈风险，因此显式栈更稳。若允许修改输入，可以把访问过的 \code{1} 改成 \code{0}；若不允许修改，就使用单独访问数组。这个选择不是风格问题，而是题目副作用要求。

正确性来自连通块定义。一次搜索从某个未访问陆地出发，会沿合法边访问所有与它连通的陆地；搜索不会越过水或边界，所以不会把两个不同岛屿合并。外层扫描保证每个连通块至少会被第一次遇到的陆地启动一次，而访问标记保证同一连通块不会重复计数。实验覆盖全水、全陆、单个陆地、斜角相邻但不连通、细长岛、输入是否可修改以及大网格递归深度。

\noindent\textbf{LeetCode 207 课程表。}

课程表题不是普通遍历，而是有向依赖图判环。每门课是节点，先修关系是有向边。如果学习课程 \code{b} 后才能学 \code{a}，可以建边 \code{b -> a}，表示 \code{b} 完成后减少 \code{a} 的依赖。暴力思路是反复找没有先修要求的课程并删除它；这正是拓扑排序的朴素模型。优化做法用入度数组保存每门课尚未满足的依赖数量，用队列保存当前入度为零的课程。

拓扑排序的过程是：初始把所有入度为零的节点入队；每弹出一门课，就把它指向的后继课程入度减一；如果某个后继入度变成零，就入队。最后如果弹出的课程数等于总课程数，说明所有依赖都能被满足；否则剩余课程处在环里或依赖环。边方向要讲清楚，否则输出顺序会反。判能否完成时，边方向反了有时也能判环，但到了课程表 II 输出顺序就会错。

正确性来自入度语义。入度为零表示当前没有未完成先修，可以安全安排；安排它只会减少后继依赖，不会破坏其他课程合法性。若最终还有节点未输出，这些节点的入度无法降到零，说明它们之间存在循环依赖。实验要覆盖无依赖、单链依赖、简单环、多个连通分量、自环、重复边和边方向检查。DFS 三色标记也能做，但需要区分“当前路径中”和“已经完成”的节点。

\noindent\textbf{LeetCode 300 最长递增子序列。}

先写二次 DP。令 \code{dp[i]} 表示必须以 \code{nums[i]} 结尾的最长递增子序列长度。对每个 \code{i}，枚举所有 \code{j < i}，若 \code{nums[j] < nums[i]}，就可以从 \code{j} 转移到 \code{i}。这个状态非常清楚，复杂度是 \complexity{O(n^2)}。它的问题是每个位置都要扫描所有历史前驱，很多历史状态只是在比较“能否成为某个长度的更小结尾”。

二分优化维护 \code{tails[len]}：长度为 \code{len + 1} 的递增子序列，当前能达到的最小结尾值。这个数组不是某条真实子序列，而是未来扩展潜力表。扫描新数 \code{x} 时，用 \code{lower_bound} 找第一个大于等于 \code{x} 的位置并替换；如果位置在末尾，就把 \code{x} 追加，说明能形成更长序列。为什么替换安全？同样长度下，结尾越小，未来越容易接上更大的数，不会比原结尾差。

正确性要强调两个层次。第一，\code{tails} 的长度始终是已经扫描前缀中可形成的某个递增子序列长度；第二，\code{tails} 每个位置保存该长度的最小可能结尾，因此替换不会减少当前答案，只会改善未来扩展空间。重复元素时使用 \code{lower_bound} 保持严格递增；若题目改成非递减，就要用 \code{upper_bound}。若要还原路径，单靠 \code{tails} 不够，需要前驱数组和位置记录。

\noindent\textbf{LeetCode 322 零钱兑换。}

这题要求凑出金额的最少硬币数，硬币可以无限使用。暴力搜索会枚举每种硬币使用次数或递归尝试所有选择，很多剩余金额会重复出现。状态可以定义为 \code{dp[value]}：凑出金额 \code{value} 的最少硬币数。初始化 \code{dp[0] = 0}，其他为无穷。对每个金额，枚举硬币，如果 \code{value >= coin} 且 \code{dp[value - coin]} 可达，就尝试更新 \code{dp[value]}。

这里要区分“最少数量”和“组合数量”。零钱兑换 I 求最少硬币数，内外循环顺序有多种写法，只要转移语义正确即可；零钱兑换 II 求组合数，循环顺序决定是否把不同顺序算作不同方案。哨兵也要谨慎：如果无穷值设得太大，\code{dp[value - coin] + 1} 可能溢出。常见做法是把无穷设为 \code{amount + 1}，因为最坏情况下使用面额一的硬币也不会超过这个数量。

正确性来自最后一枚硬币。任意最优方案最后使用某枚 \code{coin}，去掉它后剩余金额是 \code{value - coin} 的最优方案；否则如果剩余部分不是最优，就能替换得到更少硬币。实验覆盖金额为零、无法凑出、硬币面额大于目标、重复硬币、存在面额一和不存在面额一。复杂度是 \complexity{O(amount \cdot coins)}，这是伪多项式复杂度，依赖金额数值而不是输入字符长度。

\noindent\textbf{LeetCode 560 和为 K 的子数组。}

暴力枚举所有左右端点并求和，若每次重新求和是 \complexity{O(n^3)}，用前缀和可以降到 \complexity{O(n^2)}。继续优化的关键是代数关系：区间 \code{(i,j]} 的和等于 \code{prefix[j] - prefix[i]}。当扫描到当前前缀 \code{prefix[j]} 时，需要寻找多少个历史前缀等于 \code{prefix[j] - k}。这就是哈希表保存历史前缀频次的理由。

先查询还是先更新是本题核心边界。我们要找的是当前下标之前的历史前缀，所以应该先用当前前缀查询 \code{prefix - k}，再把当前前缀加入表。初始表里要放 \code{0 -> 1}，表示从数组开头到当前下标的区间也可能满足条件。若漏掉这个初始值，输入 \code{[k]} 会错误返回零。

这题能处理负数，是它和正数滑动窗口的关键区别。滑动窗口依赖右端扩展和左端收缩带来的和单调变化；一旦有负数，窗口和可能来回变化，不能安全排除左端。前缀和加哈希不需要单调性，只需要代数等式。实验覆盖负数、零、目标为零、重复前缀、从下标零开始的答案和答案数量很多。若题目改成最长区间，哈希表应保存最早位置；若改成是否存在，保存出现与否即可。

\noindent\textbf{LeetCode 743 网络延迟时间。}

题目给有向带权图，要求从起点到所有节点的最短路最大值。暴力枚举路径会爆炸；普通 BFS 只适用于等权边，因为它按边数层次扩展，而题目边权可能不同。边权非负时，Dijkstra 是合适选择：维护每个节点当前已知最短距离，用小顶堆每次取距离最小的候选扩展。

C++ 的 \code{priority_queue} 不支持 decrease-key，所以常用“压入新状态、弹出时跳过过期状态”的写法。每次松弛边 \code{u -> v}，如果 \code{dist[u] + w < dist[v]}，就更新 \code{dist[v]} 并把新状态放入堆。堆里可能还保留旧距离；弹出时如果堆顶距离不等于当前 \code{dist[node]}，说明它已经过期，直接跳过。这个懒处理是容器限制下的工程实现，不是算法偷懒。

正确性依赖非负边。每次弹出的未过期节点拥有当前全图最小的临时距离，由于所有边权非负，从其他未确定节点绕回来不可能得到更短路径，因此它的距离可以确定。实验要覆盖不可达节点、重边、起点到自己、零权边、多条等长路径和堆中过期状态。最终若有节点距离仍为无穷，返回 -1；否则返回最大最短距离。

\noindent\textbf{LeetCode 862 和至少为 K 的最短子数组。}

如果数组全是正数，长度最小且和至少为 K 可以用滑动窗口；但本题允许负数，窗口和不再单调。一个负数进入窗口可能让和变小，一个负数离开窗口可能让和变大，所以单纯移动左右端无法安全排除候选。先写前缀和：需要找 \code{prefix[j] - prefix[i] >= k} 且 \code{j - i} 最短。对每个 \code{j}，我们希望历史前缀尽量小且下标尽量靠近。

单调队列维护候选前缀下标，队列中的前缀和严格递增。处理当前 \code{j} 时，若队首和当前前缀差至少 K，就可以更新答案并弹出队首，因为这个队首作为左端已经得到最短的当前右端，未来右端更远只会长度更长。然后维护队尾：如果队尾前缀和大于等于当前前缀，且当前下标更靠右，那么队尾被当前下标支配，未来用当前下标只会得到更大的区间和和更短长度。

正确性核心是支配关系。前缀和更大且下标更早的候选不如前缀和更小且下标更晚的候选；满足条件的队首在当前右端下已经结算完毕。实验要覆盖负数破坏普通窗口、答案从开头开始、无答案、多个候选同时满足、队尾连续弹出和 K 为负或零的边界。这个题是单调队列最值得精读的一题，因为它把前缀和、最短长度和支配关系放到一起。

\noindent\textbf{LeetCode 1631 最小体力消耗路径。}

这题的路径代价不是边权和，而是路径上最大高度差。暴力可以枚举路径，但路径数量巨大。第一种高质量建模是 Dijkstra 变体：节点是格子，边权是相邻格子高度差，路径代价是当前路径最大边权。从一个格子扩展到邻居时，新代价是 \code{max(current_cost, edge_cost)}。这个代价随着路径延长不会下降，因此仍然满足按当前最小代价优先扩展的单调性。

第二种建模是二分答案。猜一个体力上限 \code{limit}，只允许走高度差不超过 \code{limit} 的边，检查左上角是否能到右下角。上限越大，可用边越多，可达性单调，所以可以二分最小可行上限。第三种建模是并查集：把所有边按高度差排序，从小到大加入，直到起点和终点连通，此时加入的边权就是答案。

三种方法利用的是不同结构。Dijkstra 把路径代价作为动态极值扩展；二分 BFS 利用可行性单调；并查集利用从小到大开放边的连通性。面试中能比较三者，比只写一种代码更有说服力。实验覆盖单格、平坦网格、必须绕路、局部高度差很大但可绕开、多个同代价路径和大网格性能。注意不要把路径代价误写成边权和，否则就变成了另一道最短路题。

\subsection{容器实验和源码对照}

\noindent\textbf{C++ 容器底层实验。}

\noindent\textbf{实验一：\code{vector} 的容量和地址。}

很多刷题代码默认 \code{vector} 是“无限数组”，但工程上必须理解它的容量模型。写一个实验：从空 \code{vector<int>} 开始连续 \code{push_back}，每次当 \code{capacity} 改变时打印旧容量、新容量和 \code{data()} 地址。你会看到容量不是每次加一，而是按某种增长策略跳变；地址也可能变化。这解释了为什么保存元素指针、引用或迭代器后继续插入是危险的。

第二步比较 \code{reserve} 和 \code{resize}。\code{reserve(1000)} 只保证容量，不改变 \code{size}，不能直接访问 \code{v[999]}；\code{resize(1000)} 会构造一千个元素，\code{size} 变为一千。很多同学为了避免扩容使用 \code{resize}，结果把未写入元素也当成有效数据，答案被默认零污染。高质量题解要说明自己需要的是容量还是大小。

第三步比较顺序扫描 \code{vector} 和 \code{list}。两者理论上都是线性，但 \code{vector} 连续内存通常有更好的缓存局部性，\code{list} 每个节点分散分配，指针追踪成本高。这个实验能纠正一个常见误解：链表插入删除复杂度好，不代表所有场景都快。若题目主要是顺序扫描和随机访问，\code{vector} 往往是更好的默认选择。

\noindent\textbf{实验二：\code{unordered_map} 的查询语义。}

哈希表题最常见错误不是哈希函数，而是 API 语义。写三版前缀和计数：第一版用 \code{operator[]} 查询，第二版用 \code{find}，第三版用 \code{contains} 后再取值。观察不存在 key 时，\code{operator[]} 会插入默认值，导致表大小变化。计数累加时这很方便，但只想判断存在时可能污染状态，也会让调试输出误以为某些前缀已经真实出现过。

再做 \code{reserve} 实验。向 \code{unordered_map} 插入大量元素，记录桶数量和负载因子变化。没有 \code{reserve} 时，插入过程中可能多次 rehash，迭代器失效，性能也会波动。刷题中不一定必须 \code{reserve}，但当输入规模很大、哈希表是热路径时，主动预留容量可以降低常数成本。需要强调的是，\code{reserve} 不能修复错误 key 设计；如果 key 构造昂贵或容易冲突，复杂度仍然会差。

最后检查 key 设计。异位词分组若把计数直接拼成字符串而不加分隔符，\code{1,11} 和 \code{11,1} 可能混淆；二维坐标若用 \code{x * M + y} 编码，要确认不会溢出且 M 选择正确；自定义结构作为 key 时，哈希函数和相等比较必须表达同一个等价关系。哈希表只是查询工具，真正决定正确性的是 key 的语义。

\noindent\textbf{实验三：\code{priority_queue} 和过期状态。}

\code{priority_queue} 只支持访问和弹出堆顶，不支持删除任意元素，也不支持 decrease-key。Dijkstra、滑动窗口中位数、任务调度和 Top K 题都必须面对这个限制。写一个 Dijkstra 小实验：每当发现更短距离，就把新状态压入堆；弹出时如果距离和 \code{dist} 数组不一致，就统计一次过期弹出。这个实验能让你直观看到“堆中有旧状态”并不影响正确性，只要使用前检查。

再写滑动窗口最大值或中位数的懒删除实验。窗口滑动时，一些元素已经过期，但它们可能不在堆顶，无法立即删除。可以用哈希表记录待删除次数，等它们浮到堆顶时再真正弹出。这个思想和 RCU 的延迟释放相似但目的不同：堆懒删除是容器能力不足下的工程技巧，RCU 是并发读者安全要求。

比较器也要单独实验。C++ 默认 \code{priority_queue<int>} 是大顶堆；小顶堆要写 \code{greater<int>}。如果存的是结构体，比较器含义是“谁的优先级更低”，很多人会写反。若比较器对相等元素处理不满足严格弱序，排序和堆行为都可能不稳定。任何堆题都要先在纸上写出堆顶应该是谁，再写比较器。

\noindent\textbf{实验四：\code{map}、\code{set} 和动态区间。}

有序容器适合前驱、后继和范围边界。写一个“我的日程安排”实验：用 \code{map<int,int>} 保存不重叠区间，key 是起点，value 是终点。插入新区间 \code{[l,r)} 时，先找第一个起点不小于 \code{l} 的后继，再检查前驱和后继是否与新区间重叠。只要这两个邻居合法，其他区间就不可能冲突，因为已有区间按起点有序且互不重叠。

这个实验能解释为什么哈希表不适合动态区间。哈希表能查某个起点是否存在，却不能快速找到离 \code{l} 最近的前驱和后继。排序数组适合离线处理，但动态插入会移动大量元素。红黑树类结构牺牲一些常数，换来稳定的对数级顺序操作。Linux 的 rbtree 和 Maple Tree 也服务于类似需求，只是多了侵入式节点、并发和内存管理约束。

实验报告要记录闭区间和半开区间的差异。若使用 \code{[l,r)}，前一个区间终点等于新区间起点不算重叠；若使用闭区间，端点相等可能算重叠。很多区间题错误都来自没定义端点语义。写比较器时也要小心相同起点：覆盖类题常需要起点升序、终点降序；合并类题只按起点升序通常足够。

\noindent\textbf{Linux 源码数据结构对照。}

\noindent\textbf{\code{list_head} 和 LeetCode 146。}

Linux 的 \code{list_head} 和力扣 LRU 的双向链表有共同点：都需要已知节点常数时间删除、移动到头部和从尾部淘汰。区别在所有权。力扣常用 \code{std::list<Node>} 让容器拥有节点，哈希表保存迭代器；内核常把 \code{list_head} 嵌进业务对象，链表只保存链接关系，不负责对象分配和释放。这个侵入式设计让同一个对象可以同时挂进多个链表，例如一个页面对象可以进入 LRU 链表，也可以进入其他状态链表。

源码阅读时不要先追每个宏，而要先画对象布局：业务对象里有哪些 \code{list_head} 字段，每个字段对应哪条链，插入和删除由谁保护，删除后对象是否还能被其他索引找到。刷题中的 LRU 删除尾部只要维护哈希表和链表；内核里删除一个对象可能还要处理引用计数、锁、RCU 或延迟释放。相同点是“节点身份必须稳定”，不同点是“生命周期复杂得多”。

实验可以写两个版本的 LRU。第一版使用 \code{std::list} 和 \code{unordered_map}；第二版让 \code{Entry} 自己带 \code{prev}、\code{next} 指针，哈希表保存 \code{Entry*}。比较两版的删除顺序：先从链表摘除，再从哈希表删除，最后释放对象。如果先释放对象，链表或哈希表就会留下悬空指针。这个实验能把链表题、缓存题和内核侵入式结构连起来。

\noindent\textbf{\code{hlist} 和哈希冲突。}

哈希表的桶内冲突可以用链表处理。Linux 的 \code{hlist} 面向大量桶头做了空间优化：桶头较轻，节点保存足够信息以支持已知节点删除。刷题时 \code{unordered_map} 把桶数组、负载因子和冲突链隐藏起来，但理解桶结构能帮助你解释为什么哈希表平均常数不是无条件保证。

对照题包括两数之和、字母异位词分组、和为 K 的子数组、账户合并。每道题都不是“用了哈希就完了”，而是要回答 key 是什么、冲突时如何确认相等、value 保存次数还是下标、什么时候插入、什么时候删除。内核里这些问题更具体：对象节点是否已在桶中，删除时是否有读者正在遍历，key 对应的对象生命周期由谁保证。

实验可以写一个固定 16 个桶的小哈希表。第一版桶内使用普通单链表，删除已知节点时必须从桶头找前驱；第二版让节点保存前驱链接位置，删除时直接改指针。这个实验能解释 hlist 的设计动机：不是为了让算法更神秘，而是为了在桶很多、删除频繁的场景下减少空间和扫描成本。

\noindent\textbf{rbtree 和动态有序集合。}

Linux rbtree 与 C++ \code{map} 都是动态有序结构，但接口风格不同。C++ 容器把节点、比较器和内存管理封装起来；Linux rbtree 通常把树节点嵌入业务对象，比较逻辑由调用者在插入和查找时写出。这样做适合内核对象按不同字段进入不同索引，也避免通用容器隐藏分配和生命周期。

对照刷题时，可以把 rbtree 看作“需要前驱后继的动态集合”的底层代表。日程安排、区间冲突、滑动窗口有序统计、数据流排名，都不是哈希表擅长的问题。红黑树的旋转不会改变中序顺序，只是在保持有序语义的前提下修复平衡。理解这一点后，再看 \code{lower_bound}、前驱、后继和区间插入，会更容易说清为什么只检查相邻区间。

源码阅读建议选一条简单路径：从查找插入位置开始，看调用者如何比较 key；再看链接新节点；最后看插入修复如何通过旋转和变色维持平衡。不要第一天就背红黑树五条性质的所有修复分支。刷题阶段也一样：先证明有序性和相邻关系，再谈容器实现。

\noindent\textbf{XArray、Trie 和稀疏整数索引。}

XArray 面向整数索引到指针的映射，适合稀疏 ID 空间和内核对象管理。它和 Trie 有相似思想：把 key 分段，沿层级向下查找，只为实际存在的路径分配节点。与普通哈希相比，它更重视按索引组织、范围推进、标记和并发语义；与巨大数组相比，它避免为空洞分配大量空间。

刷题中可以用位 Trie、前缀树、坐标压缩来建立直觉。最大异或题把整数按二进制位分层；单词搜索 II 把字符串按字符分层；坐标压缩把稀疏大值域映射到紧凑下标。XArray 的工程环境更复杂，但核心问题仍是：key 空间很大，实际对象稀疏，等值查找之外还可能需要按序扫描和状态标记。

实验可以写一个简化的二进制 Trie 保存整数集合，再支持查找某个整数是否存在、查找不小于某个值的下一个元素。这个实验不会复刻 XArray，但能训练分层索引的感觉。报告中要比较数组、哈希表、平衡树和 Trie：它们分别适合连续小值域、稀疏等值查找、动态有序查找、可分段 key。

\noindent\textbf{Maple Tree 和动态区间。}

Maple Tree 服务于范围映射场景，典型直觉是虚拟地址区间管理。区间问题比单点 key 更难：查询一个地址需要找到覆盖它的区间；插入新区间需要检查前后邻居是否重叠；删除可能删除整段，也可能把原区间切成两段。力扣 56 和 57 是离线区间题，Maple Tree 面对的是动态范围集合和高频查询。

刷题阶段可以先用 \code{std::map} 写简化区间表。key 是区间起点，value 是终点和属性。点查询时找起点不大于地址的最后一个区间，再判断终点是否覆盖；插入时检查前驱和后继；删除中间一段时要拆分。这个小实验能帮助理解 Maple Tree 的问题动机，而不是把源码当成神秘结构背诵。

工程差异在于读写并发、内存分配和查询路径。一个简单 \code{map} 版本可以帮助学习区间不变量，却不能替代内核结构。源码阅读时要问：节点里保存哪些范围信息，读路径如何快速找到覆盖区间，更新如何保持相邻关系，旧节点什么时候释放。这些问题比记字段名更稳定。

\noindent\textbf{伙伴系统、SLUB 和分配成本。}

算法题通常把内存分配当成免费，但 Linux 内存管理会把分配成本放到中心。伙伴系统按 2 的幂次阶管理连续页，分配时从合适阶取块，不够就拆更高阶；释放时检查伙伴是否空闲，能合并就合并。它强调连续性和对齐，不只是总空闲量。一个系统可能总空闲页很多，却没有足够大的连续块。

SLUB 解决小对象频繁分配问题。同类型对象从缓存中分配，避免每次都走通用页分配路径，也提升局部性。回到刷题，链表节点、哈希节点、树节点都可能带来大量小对象分配；这就是为什么连续数组有时远快于节点容器。复杂度分析给出数量级，分配器和缓存局部性决定常数。

实验可以写简化伙伴系统：管理 \code{2^10} 个页框，维护 0 到 10 阶空闲链表。分配某阶时逐级拆分，释放时按伙伴编号尝试合并。再写一个对象池，为固定大小节点重复分配和释放，比较它与频繁 \code{new} 的差异。这个实验能把位运算、链表、区间合并和工程性能放在同一张图里。

\subsection{错误用例、复盘和训练路线}

\noindent\textbf{错误用例库。}

\noindent\textbf{数组和窗口。}

数组题的错误用例要专门打破区间语义。删除重复项要测空数组、全重复、无重复和只保留两次的变体；颜色分类要测全二、全零、逆序和交换后二次检查；接雨水要测单调递增、单调递减和平台高度。窗口题则要测试左端不能回退的情况，例如无重复子串的 \code{abba}；还要测试目标字符重复，例如最小覆盖子串中目标含两个相同字符。

设计反例时，不要只看最终答案，要让错误状态暴露出来。滑动窗口题可以打印左端、右端、计数和答案；双指针题可以打印每次为什么移动某一侧。一个好的反例能说明错误发生在初始化、循环条件、状态更新还是答案读取，而不是只告诉你“答案错了”。这也是把题解变成教材的关键。

\noindent\textbf{哈希和前缀和。}

哈希题的错误用例集中在查询和更新顺序、key 设计、保存次数还是下标。两数之和要测补数等于当前值；和为 K 的子数组要测从下标零开始的答案和目标为零；连续子数组和要测长度至少为二；连续数组要测重复前缀且要求最长。异位词分组要测空字符串、重复单词和计数 key 歧义。

前缀和题尤其适合对拍。写一个平方暴力版，再随机生成包含负数、零和重复值的数组，与哈希版比较。很多窗口解法在正数数组上正确，一加负数就失败；对拍能强迫你说明算法依赖的是单调性还是代数等式。若目标从“计数”改成“最长”，哈希表 value 也要从频次改成最早下标，这类变化必须进反例库。

\noindent\textbf{二分和答案空间。}

二分错误不是只差一，而是区间语义没有固定。测试要覆盖空数组、长度一、长度二、目标小于所有元素、目标大于所有元素、目标重复出现、答案在边界。答案二分还要测试刚好可行和刚好不可行的相邻值。每个测试都打印 \code{left}、\code{mid}、\code{right} 和谓词值，检查被排除的一半是否真的不含答案。

对于船运包裹、爱吃香蕉、分割数组这类题，检查函数本身就是贪心题。反例要针对检查函数，而不是只针对外层二分。比如容量小于最大包裹必须不可行；天数为一时容量必须至少是总和；天数等于包裹数时容量下界就是最大包裹。若检查函数没有这些边界，二分循环再标准也会返回错误答案。

\noindent\textbf{栈、堆和单调结构。}

单调栈要测试严格递增、严格递减、相等元素和哨兵。柱状图最大矩形中，递增数组能暴露末尾未结算问题；递减数组能暴露宽度计算问题；全相等能暴露相等元素处理策略。每日温度要测试没有更高温、相同温度不算更高、最后一个元素答案为零。

堆题要测试过期状态和比较器方向。Top K 高频要测试频次相同和 K 等于不同元素数；数据流中位数要测试偶数个元素平均溢出；Dijkstra 要测试重边和旧堆状态；滑动窗口中位数要测试要删除的元素不在堆顶。每个堆实验都应该打印堆顶是否真实可用，而不是默认堆顶一定有效。

\noindent\textbf{树、图和并查集。}

树题要测试空树、单节点、链式树、完全树、重复值、全负值和极端深度。验证 BST 的反例必须包含“局部父子合法但全局不合法”，例如某个右子树深处节点小于根。最大路径和要测试全负数，答案不能初始化为零。构造树题要测试节点值唯一性，如果值不唯一，单靠值到中序下标的哈希表就不成立。

图题要测试自环、重边、不连通、多个源点、状态资源不同但位置相同。打开转盘锁要测试起点死亡和目标就是起点；障碍消除最短路要测试同一格子剩余消除次数不同；课程表要测试自环和两个节点互相依赖。并查集要测试重复合并、已经连通的边、字符串编号和合并失败代表环，而不是路径本身。

\noindent\textbf{动态规划。}

DP 反例要打状态定义。最大子数组和要测试全负数；不同路径要测试一行一列；障碍物版本要测试起点或终点障碍；编辑距离要测试空串；背包要测试容量为零、总和奇偶、一个物品被重复使用的循环方向错误；零钱兑换 II 要测试组合和排列的差异。空间压缩题还要测试更新顺序，例如 0-1 背包正序会错误重复使用当前物品。

打印 DP 表是最有效的实验。编辑距离表能看到插入、删除、替换；LCS 表能看到不连续子序列和子串的区别；股票状态机能看到持有、卖出、休息之间的约束。若一个状态无法解释表中某个格子的来源，就说明状态定义还不够稳定。不要急着压缩空间，先让完整状态表正确。

\noindent\textbf{复盘报告模板。}

每道题复盘写一页即可，但必须有固定信息。第一行写题号和题名，例如 LeetCode 560 和为 K 的子数组。第二行写题面模型：对象、关系、输出。第三行写暴力基线：如何保证不漏候选，复杂度是多少。第四行写瓶颈：哪一步重复扫描、重复计算或无法快速查询。第五行写优化结构：保存什么，查询什么，更新什么，删除什么。第六行写正确性：覆盖和排除各一句。第七行写边界用例：至少三个能打破错误实现的输入。第八行写 C++ 注意点：容器、溢出、迭代器、所有权或副作用。

这份报告不需要文学化，越具体越好。比如“用哈希优化”不是合格描述，“哈希表保存当前下标之前每个前缀和出现次数，扫描到前缀 \code{s} 时先查询 \code{s-k} 再累加当前前缀”才是合格描述。比如“注意边界”也不是合格描述，“初始放入 \code{0 -> 1}，否则从下标零开始的区间会漏掉”才是合格描述。

最后要有实验记录。把暴力版和优化版同时保留，用随机小数据对拍；把边界用例写成单元测试；把大规模输入用于观察复杂度增长。对拍不是浪费时间，它能把“我觉得对”变成“我知道错在哪里”。当你能用一页报告解释一道题，再用相似报告解释变体，刷题才真正转化为面试能力和工程能力。

\noindent\textbf{高阶经典题补充推导。}

\noindent\textbf{LeetCode 4 寻找两个正序数组的中位数。}

这题经常被当成二分难题背答案，其实先从暴力看更清楚。两个数组已经有序，最直接的方法是归并到一半位置，复杂度 \complexity{O(m+n)} 或 \complexity{O((m+n)/2)}。它没有利用“只需要中位数而不是完整合并结果”这个条件。中位数的本质是把全部元素切成左右两半，左半元素数量等于右半或多一个，并且左半最大值不大于右半最小值。

优化做法是在较短数组上二分切分位置。假设第一个数组左边取 \code{i} 个元素，第二个数组左边就必须取 \code{half - i} 个元素。这样数量条件自动满足，只剩大小条件：第一个数组左侧最大不大于第二个数组右侧最小，第二个数组左侧最大不大于第一个数组右侧最小。如果第一个数组左侧太大，说明 \code{i} 取多了，要向左；如果第二个数组左侧太大，说明 \code{i} 取少了，要向右。

边界处理依赖哨兵。切分在数组开头时，左侧最大可视为负无穷；切分在数组末尾时，右侧最小可视为正无穷。这样不需要为每个边界写特殊分支。正确性来自切分条件：数量正确且左右有序时，中位数只可能由左侧最大和右侧最小决定。实验要覆盖一个数组为空、两个数组长度差很大、总长度奇偶、重复元素、所有元素都在一个数组左侧或右侧。讲这题时重点不是背代码，而是说明二分的对象是切分位置，不是某个具体值。

\noindent\textbf{LeetCode 10 正则表达式匹配。}

这题的困难在于星号。暴力递归从字符串下标和模式下标出发，遇到普通字符或点号就匹配一个字符，遇到星号就尝试匹配零次、一次、两次甚至更多次。暴力会在同一对下标上反复尝试，复杂度爆炸。DP 状态可以定义为 \code{dp[i][j]}：字符串前 \code{i} 个字符能否匹配模式前 \code{j} 个字符。这个定义把递归中的下标对固定下来，避免重复搜索。

转移要围绕模式最后一个字符。如果 \code{p[j-1]} 不是星号，那么只有当前字符能匹配时，才能从 \code{dp[i-1][j-1]} 转移。若 \code{p[j-1]} 是星号，它修饰的是 \code{p[j-2]}。星号匹配零次时，看 \code{dp[i][j-2]}；星号匹配一次或多次时，必须当前字符能被 \code{p[j-2]} 匹配，然后看 \code{dp[i-1][j]}，表示消耗一个字符串字符后仍由这个星号继续处理。

初始化是最容易漏的地方。空字符串可以被 \code{a*}、\code{a*b*} 这类模式匹配，所以第一行需要按星号零次匹配向前传播。正确性来自模式末尾分类：任意合法匹配的最后一步要么消耗一个普通模式字符，要么让星号匹配零次，要么让星号至少匹配一次，这三类覆盖全部可能。实验覆盖空串、只含星号组合的模式、点号、星号连续修饰的结构、完全不匹配和长重复字符串。

\noindent\textbf{LeetCode 25 K 个一组翻转链表。}

这题不是单纯反转链表，而是“先确定边界，再局部反转，再接回”。暴力直觉可以把链表节点放进数组，每 K 个一组反转数组片段，再重连节点。但原地链表解法要训练的是指针边界。每一轮从上一组尾部出发，向后找第 K 个节点；如果不足 K 个，剩余节点保持原样。找到组尾后，先保存下一组头，否则反转过程中会丢失后续链表。

指针关系可以固定成四个名字：\code{group_prev} 是上一组尾部，\code{group_head} 是当前组头，\code{group_tail} 是当前组尾，\code{next_group} 是下一组头。反转当前组后，原来的 \code{group_tail} 成为新头，原来的 \code{group_head} 成为新尾。接回时，\code{group_prev->next} 指向新头，新尾的 \code{next} 指向 \code{next_group}，然后 \code{group_prev} 移到新尾。

正确性来自链表结构不变量：已经处理的前缀是按 K 分组翻转后的合法链；当前组在反转前已经确认长度足够；未处理后缀仍然由 \code{next_group} 保存且可达。实验要覆盖空链表、长度小于 K、长度等于 K、长度不是 K 倍数、K 为一、连续多组和删除头部附近边界。C++ 中不要随意释放输入节点，除非题目明确要求；这里重连节点即可。

\noindent\textbf{LeetCode 32 最长有效括号。}

暴力方法枚举每个子串并检查括号是否合法，复杂度很高。栈解法的核心是保存边界。栈里放尚未匹配的左括号下标，同时保留最后一个无法匹配的右括号位置。一个常见写法是在栈中先放 \code{-1} 作为哨兵：遇到左括号入栈；遇到右括号先弹出一个左括号位置，如果弹完后栈空，说明当前右括号无法匹配，把当前下标作为新边界入栈；否则当前有效长度就是 \code{i - stack.top()}。

为什么这个长度公式正确？栈顶保存的是当前有效后缀左侧最近的非法边界，或者尚未匹配的左括号位置。当前右括号匹配成功后，从栈顶之后到当前下标之间都是一个有效括号串，因此长度是差值。如果没有哨兵，边界会被很多特殊分支打散。DP 解法也可做，\code{dp[i]} 表示以 \code{i} 结尾的最长有效括号长度，但转移需要回看匹配左括号位置，初学阶段栈更直观。

实验要覆盖空串、全左括号、全右括号、简单嵌套、相邻有效段、非法右括号打断后重新开始。调试时打印栈内容和当前答案，特别观察 \code{)()())} 这类样例。正确性不是“栈能匹配括号”这么简单，而是栈顶始终保存当前有效后缀左侧的边界。

\noindent\textbf{LeetCode 41 缺失的第一个正数。}

排序能做，哈希集合也能做，但题目要求线性时间和常数额外空间。关键观察是：长度为 n 的数组，答案只可能在 \code{1..n+1} 之间。小于一和大于 n 的数不会影响前 n 个正整数是否存在；每个值 \code{x} 如果在 \code{1..n} 内，它的理想位置是下标 \code{x-1}。于是可以把数组本身当作哈希表，通过交换把每个合法值放到自己的位置。

循环时，对每个下标 \code{i}，只要 \code{nums[i]} 在 \code{1..n} 内，并且目标位置上的值不是它，就交换 \code{nums[i]} 和 \code{nums[nums[i]-1]}。这里必须检查目标位置是否已经有相同值，否则重复数字会导致无限交换。所有能放到正确位置的数处理完后，再从左到右找第一个 \code{nums[i] != i+1} 的位置，答案就是 \code{i+1}；如果都正确，答案是 \code{n+1}。

正确性来自位置映射。任何在 \code{1..n} 内出现过的值，只要经过交换循环，就会被移动到对应位置，除非该位置已经有相同值。无关值被留在某些位置，不影响其他合法值继续归位。实验覆盖空数组、全负数、含零、重复正数、已经连续、缺一、答案为 n+1 和极端乱序。这个题训练的是原地哈希，不是普通排序。

\noindent\textbf{LeetCode 51 N 皇后。}

暴力可以在棋盘每个格子上选择放或不放皇后，但这会产生巨大搜索树。第一层优化是逐行放置，因为合法解中每行必须恰好一个皇后；这样每层只需要选择列。再进一步，用三个集合判断冲突：列集合、主对角线集合和副对角线集合。对于位置 \code{row,col}，主对角线可用 \code{row-col} 标识，副对角线可用 \code{row+col} 标识。

回溯状态包括当前行号、已经放置的列、两类对角线和棋盘路径。进入某个选择时，把对应列和对角线加入集合；递归到下一行；返回时撤销。剪枝正确性来自皇后的攻击规则完全由同列和两类对角线覆盖。不要用“遍历已有皇后逐个检查”作为最终写法，它虽然正确，但每次检查重复扫描；集合把冲突判断降成常数时间。

实验可以统计搜索节点数。先写无剪枝版本或只检查列的版本，再加入对角线集合，比较 n 增大时节点数差异。边界包括 n 为一、n 为二和三无解、n 为四有两个解。输出棋盘时要注意路径拷贝成本：每次找到答案再构造字符串，或维护可修改棋盘并在答案处复制。这个题的重点是搜索树建模和剪枝证明。

\noindent\textbf{LeetCode 72 编辑距离。}

编辑距离是二维 DP 的代表。暴力递归从两个字符串末尾开始，如果末尾相同，就处理两个更短前缀；如果不同，就尝试删除、插入、替换三种操作。不同路径会反复处理同一对前缀长度，于是定义 \code{dp[i][j]}：把 \code{word1} 前 \code{i} 个字符变成 \code{word2} 前 \code{j} 个字符的最少操作数。

初始化来自空串。把前 \code{i} 个字符变成空串，需要删除 \code{i} 次；把空串变成前 \code{j} 个字符，需要插入 \code{j} 次。转移时，如果两个末尾字符相同，\code{dp[i][j] = dp[i-1][j-1]}；否则三种候选分别是删除 \code{word1} 末尾、向 \code{word1} 插入 \code{word2} 末尾、替换末尾，对应 \code{dp[i-1][j] + 1}、\code{dp[i][j-1] + 1}、\code{dp[i-1][j-1] + 1}。

正确性来自最后一步分类。任意最优编辑序列的最后一步必然是三种操作之一，或者末尾字符相同不需要操作；每种操作前的子问题都由更短前缀定义。实验要打印 DP 表，手动解释每个格子来源。若操作代价不同，转移权重要改；若只允许插入删除，则问题和最长公共子序列相关。空间压缩可以做，但初学时先写二维表，避免左上角旧值被覆盖。

\noindent\textbf{LeetCode 212 单词搜索 II。}

如果对每个单词单独在棋盘上 DFS，会重复搜索大量相同前缀。优化来自 Trie：把所有单词的前缀合并成一棵树，网格 DFS 时同步沿 Trie 前进。如果当前路径在 Trie 中已经没有对应子节点，就可以立即剪枝；如果 Trie 节点标记了完整单词，就把它加入答案。这样搜索空间由“单词数乘以棋盘搜索”变成“棋盘路径和词典前缀共同约束”。

状态包括网格位置、当前 Trie 节点和访问标记。进入一个格子前，先看当前 Trie 节点是否有该字符子节点；没有就返回。有则进入子节点，标记格子已访问，向四周扩展，退出时恢复访问标记。找到单词后，为避免重复加入，可以把节点上的单词标记清空；进一步优化还可以在子树为空时剪掉 Trie 叶子，但要小心节点生命周期。

正确性来自前缀共享。任意答案单词对应 Trie 中从根到终止节点的一条路径，网格 DFS 会枚举所有不重复格子的相邻路径；只有当前缀存在时才继续，不会漏掉可能单词。实验覆盖重复单词、同一单词多条路径、棋盘字符频次不足、单字符单词、访问标记恢复和 Trie 剪枝。C++ 实现中，固定小写字母可以用 26 个指针或下标数组；若字符集更大，哈希子节点更通用但常数更高。

\noindent\textbf{LeetCode 312 戳气球。}

这题最重要的转折是从“第一个戳谁”改成“最后戳谁”。如果按第一个戳的气球思考，左右邻居会随着后续操作变化，状态不稳定。若考虑某个开区间内最后被戳的气球 \code{mid}，那么在它最后被戳时，区间左右边界气球仍然存在，收益可以稳定写成 \code{nums[left] * nums[mid] * nums[right]}，左右两侧已经独立处理完。

令 \code{dp[left][right]} 表示开区间 \code{(left,right)} 内所有气球戳完的最大收益。枚举最后戳的 \code{mid}，转移为左区间收益加右区间收益加最后一次收益。为了处理边界，在数组两端补一。遍历顺序按区间长度从小到大，因为大区间依赖更短区间。这个状态定义看起来绕，但它解决了邻居变化的问题。

正确性来自最后一步分解。任意最优方案在区间内都有最后一个被戳的气球；在它之前，左右两侧气球的操作互不影响，因为最后时左右边界固定。枚举所有 \code{mid} 就覆盖全部最优方案可能的最后一步。实验覆盖一个气球、两个气球、包含一、较大值在中间、不同最后戳选择。这个题是区间 DP 的核心样本：只要正向操作让状态变化复杂，就尝试倒过来看最后一步。

\noindent\textbf{LeetCode 460 LFU 缓存。}

LFU 比 LRU 多一个频次维度。题目要求淘汰访问频率最低的 key；若频率相同，淘汰最久未使用的 key。暴力可以每次淘汰时扫描所有 key 找最低频和最旧时间，复杂度太高。要做到常数时间，需要两个索引：key 到节点信息的哈希表，频次到双向链表的哈希表。每个频次链表内部按最近使用顺序排列。

访问一个 key 时，节点频次从 f 变成 f+1，要从旧频次链表删除，再插入新频次链表头部。如果旧频次链表为空，并且它正好是当前最小频次，就把最小频次加一。插入新 key 时，频次为一，放入频次一的链表头部，最小频次重置为一。容量满时，从最小频次链表尾部淘汰最旧节点，并从 key 哈希表删除。

正确性来自两个不变量：key 表能定位每个节点；频次表中，每个链表只保存该频次节点，且从头到尾是新到旧顺序；\code{min_freq} 始终指向当前存在节点的最低频次。实验覆盖容量为零、容量为一、重复访问同一 key、更新已有 key、频次相同按 LRU 淘汰、旧频次链表清空。这个题非常接近工程索引维护：多个结构必须同步更新，任何一步漏掉都会产生悬空节点或错误淘汰。

\noindent\textbf{Linux 专题阅读路线。}

\noindent\textbf{第一周：链表和对象身份。}

第一周只读链表，不要贪多。目标是理解 \code{list_head} 为什么嵌进对象，\code{container_of} 为什么能从成员地址回到外层对象，遍历时删除为什么需要 safe 版本。读源码时画三张图：空链表哨兵、自身在链表中的节点、删除当前节点前保存下一个节点。然后写一个小实验：同一个对象带两个链表节点，分别挂入 LRU 链和脏链。实验报告说明删除其中一条链的节点不会自动从另一条链删除。

这一周要和力扣链表题对照。反转链表训练保存后继，K 个一组翻转训练分组边界，LRU 训练已知节点移动和尾部淘汰。内核链表把这些小技巧放到真实生命周期里：对象可能被多个模块引用，删除链表节点不等于释放对象。读懂这一点，再看内核其他数据结构就不会只停留在指针操作。

\noindent\textbf{第二周：哈希链和有序树。}

第二周读 hlist 和 rbtree。hlist 关注桶头空间、冲突链和已知节点删除；rbtree 关注动态有序索引和调用者比较。读 hlist 时，写一个小哈希表，比较普通单链表删除和保存前驱链接位置的删除。读 rbtree 时，不要先背全部修复情况，而是先看调用者如何找到插入位置，再看节点如何链接到父节点，最后理解旋转保持中序顺序。

对应刷题题型是哈希与区间。两数之和、前缀和计数、异位词分组训练 key 设计；日程安排、合并区间动态版、滑动窗口有序统计训练前驱后继。你要能说出什么时候哈希表不够：只要需要“小于等于某值的最大 key”“大于等于某值的最小 key”或按顺序扫描，哈希表就不是合适结构。

\noindent\textbf{第三周：XArray、Maple Tree 和范围。}

第三周读整数索引和范围索引。XArray 可以先用 Trie 类比：把整数 key 分段，沿层级定位对象，适合稀疏空间。Maple Tree 可以先用 \code{std::map} 区间表类比：给地址找覆盖区间，插入区间检查邻居，删除区间可能拆分。源码细节很多，但第一遍只回答三个问题：key 或范围如何组织，查找路径如何走，更新后旧节点如何处理。

对应刷题题型包括 Trie、坐标压缩、区间合并、插入区间、删除被覆盖区间。区别在于内核还要处理并发和对象生命周期。不要把 Maple Tree 简化成“高级红黑树”，它面向的是范围映射和读路径效率。读完这一周，应能写出一个简化 VMA 表，支持点查询、插入不重叠区间、删除子区间，并解释它缺少哪些内核级能力。

\noindent\textbf{第四周：RCU 和内存分配。}

第四周读 RCU、伙伴系统和 SLUB。RCU 的关键词是读侧轻量、写侧复制或替换、旧对象延迟释放。把删除拆成逻辑删除和物理释放：先从索引摘除，让新读者找不到；再等旧读者离开；最后释放。这个思想能帮助你理解为什么源码里看似“已经删掉”的对象还不能马上归还内存。

伙伴系统和 SLUB 则回答分配成本。伙伴系统关心连续页和阶数，SLUB 关心固定大小对象复用。对应刷题里的对象池、链表节点分配、哈希表节点常数、数组缓存局部性。第四周实验可以写一个读多写少配置表：读者只读全局指针，写者复制新表并发布旧表到延迟列表；再写简化伙伴分配器，观察总空闲量足够但连续块不足的情况。

\noindent\textbf{专题实验清单。}

\noindent\textbf{实验一：暴力与优化对拍框架。}

为每类题保留一个暴力版。数组和窗口题用小规模枚举；图题用小图全搜索或慢 BFS；DP 题用记忆化递归和表格互相校验；哈希题用平方枚举校验。对拍框架随机生成小输入，比较暴力版和优化版输出，一旦不同就打印输入、两版结果和关键状态。这个实验能极大减少“看起来对”的错觉。

对拍不是为了替代证明。它的作用是暴露边界，证明的作用是解释为什么所有边界都被覆盖。报告里要写：随机输入如何生成，哪些约束被覆盖，暴力版为什么可信，优化版错在哪里。比如和为 K 的子数组，要让随机数组包含负数和零；滑动窗口正数题，则要分别生成全正和含负两组，观察算法适用边界。

\noindent\textbf{实验二：复杂度曲线。}

选三道题画复杂度曲线：两数之和的平方枚举和哈希版，柱状图最大矩形的暴力扩展和单调栈版，岛屿数量的重复搜索和访问标记版。每个版本运行多个输入规模，记录耗时中位数。不要把打印放进计时代码，不要只运行一次，不要在 Debug 编译下得出结论。实验结果不必精确符合数学曲线，但增长趋势应能说明优化替换了哪一步。

复杂度曲线能训练工程判断。有时 \complexity{O(n \log n)} 的排序在小数据上比复杂的线性算法更快；有时哈希表因为 key 构造成本高而慢；有时链表理论删除快，但缓存局部性差导致整体慢。教材里的复杂度是第一层判断，实验让你看到常数和内存布局。

\noindent\textbf{实验三：边界测试表。}

为每个题型建立边界表。二分题固定测试空、一、二、边界外、重复；窗口题固定测试空串、全重复、目标缺失、刚好覆盖、窗口多次收缩；图题固定测试不连通、自环、重边、多源、资源维度；DP 题固定测试空状态、全负、不可达、空间压缩顺序。每次写题先从表里选用例，不要等错了再补。

边界表的价值在于复用思维，而不是复用代码。比如“重复元素”在三数之和里是去重问题，在二分里是边界问题，在堆里是比较器 tie-breaker 问题，在哈希里是频次问题。表格要记录同一个输入特征在不同题型中的含义，这样复习时才能横向迁移。

\noindent\textbf{实验四：源码阅读报告。}

每次读源码只回答一个小问题。例如读 \code{list_head} 时，只回答删除节点改了哪几条边；读 rbtree 时，只回答插入前调用者如何比较 key；读 XArray 时，只回答索引如何分层；读 RCU 时，只回答旧对象什么时候释放。报告必须包含文件路径、接口名、对象结构、不变量、和一个刷题类比。

不要把源码报告写成字段翻译。字段名会随版本变，设计问题更稳定。一个合格报告应写出：这个结构服务的访问模式是什么，为什么普通数组、哈希表或链表不够，插入和删除维护什么关系，并发或生命周期带来什么额外约束。把这四个问题写清楚，比背十个宏更有价值。

\noindent\textbf{面试讲解脚本。}

\noindent\textbf{从题目到方案。}

高质量口述可以固定成六句。第一句复述题面模型，不加算法名；第二句给暴力解，说明它如何覆盖所有候选；第三句指出瓶颈，说明哪一步重复；第四句给关键观察，说明某个状态可以复用或某些候选可以排除；第五句讲数据结构职责，说明保存、查询、更新、删除；第六句讲正确性和复杂度。这个顺序能防止直接背答案，也能让面试官看到推导过程。

以 LeetCode 560 为例：题面模型是区间和计数；暴力枚举所有左右端点；瓶颈是每个右端要找很多历史左端；关键观察是 \code{prefix[j] - prefix[i] == k}；数据结构是哈希表保存历史前缀和频次；正确性来自任意合法区间都对应一对前缀。这个脚本比“用前缀和加哈希表”更完整。

\noindent\textbf{面对追问。}

追问通常改一个条件。数组从全正改成含负，滑动窗口可能失效；返回一个答案改成返回所有答案，输出规模和去重会变化；只判断存在改成计数，哈希表 value 要从布尔变成频次；求最短改成求最长，保存最早位置或单调队列；普通树改成 BST，可以利用中序和上下界。面对追问先问条件变了哪里，再判断原状态是否仍然足够。

不要用“这个也差不多”回答变体。三数之和和四数之和相似，但四数之和要注意外层剪枝和 \code{std::int64_t}；接雨水和盛水最多相似处很少，一个计算每个位置贡献，一个选择两条边；零钱兑换 I 和 II 都是背包，但一个求最少数量，一个求组合数，循环语义不同。能指出相似题的不同点，比会背更多题更重要。

\noindent\textbf{讲代码质量。}

面试代码要能说明工程取舍。为什么用 \code{vector} 而不是 \code{list}，为什么哈希表要先查再插入，为什么堆里保存下标而不是值，为什么距离用 \code{std::int64_t}，为什么比较器这样写，为什么递归可能有栈风险。代码短不是唯一目标，状态清楚、边界自然、复杂度可信才是高质量。

如果写类，例如 LRU、并查集、Trie，要说明成员职责。并查集的 \code{parent} 保存代表关系，\code{size} 或 \code{rank} 控制合并；LRU 的哈希表负责定位，链表负责顺序；Trie 节点负责前缀转移和单词结束。类方法内部访问成员时保持一致风格，复杂逻辑拆成小函数，如 \code{move_to_front}、\code{remove_node}、\code{link_after}。这不仅让代码好看，也让不变量有固定位置。

\noindent\textbf{最后的学习路线。}

\noindent\textbf{第一轮：题型地图。}

第一轮目标不是做难题，而是建立题型地图。数组、窗口、哈希、二分、栈、堆、树、图、回溯、DP、贪心、并查集都要各做一组代表题。每题写出暴力、瓶颈、优化结构、复杂度和两个边界。第一轮可以看题解，但看完必须用自己的话写一页复盘。不会证明也没关系，但要知道证明缺口在哪里。

\noindent\textbf{第二轮：证明和边界。}

第二轮不追求新题数量，而是补证明。二分题证明排除一半，窗口题证明左端不会漏，单调栈证明弹出时答案确定，贪心题证明局部选择不劣，DP 题证明最后一步覆盖全部可能，图题证明首次到达或最短距离。每个题型至少准备三个反例，能解释错误代码为什么会被打破。

\noindent\textbf{第三轮：工程化表达。}

第三轮把代码改成工程上也能接受的版本。统一变量语义，处理整数溢出，避免不必要复制，确认容器 API 副作用，给复杂状态写辅助函数。对照 C++ 容器实验和 Linux 源码专题，重新审视每道题的数据结构选择。最终目标是：你不仅能 AC，还能解释为什么这个结构适合这个访问模式，为什么边界不会漏，为什么复杂度可信。

\noindent\textbf{高频综合题继续精读。}

\noindent\textbf{LeetCode 85 最大矩形。}

最大矩形是柱状图最大矩形的二维化。暴力枚举左上角和右下角，再检查矩形内是否全是一，复杂度会非常高；即使用二维前缀和把检查降成常数，枚举矩形仍然是四次方级别。真正的观察是把每一行当作矩形底边：如果当前格子是 \code{1}，它上方连续一的数量可以累加成柱高；如果当前格子是 \code{0}，柱高清零。于是每一行都转化成一组柱状图高度，对这组高度运行 LeetCode 84 的单调栈算法。

这个转化有两个关键点。第一，任何全一矩形都有一个底边行，当扫描到这行时，它会在柱状图中表现为一段连续高度至少为矩形高度的柱子；单调栈会在某根柱子作为最矮高度时结算出它。第二，遇到零必须把高度清零，因为以当前行为底的矩形不能穿过零。不要把每行高度重新向上扫描，那会把复杂度又拉高；高度数组应当逐行增量维护。

实验时先在纸上画一个小矩阵，逐行写出高度数组，再运行柱状图栈。边界包括全零、全一、单行、单列、零把高度截断、最大矩形不在最后一行。C++ 中矩阵字符常是 \code{'0'} 和 \code{'1'}，不要和整数零一混淆。讲复杂度时要写 \complexity{O(mn)}，因为每一行的柱状图算法是线性的。

\noindent\textbf{LeetCode 97 交错字符串。}

交错字符串的暴力递归非常直观：目标串当前位置可以来自第一个字符串，也可以来自第二个字符串，只要字符相等就递归尝试。问题是同一对下标 \code{i,j} 会被不同路径反复到达。状态定义应当保留两个来源串各用了多少字符：\code{dp[i][j]} 表示 \code{s1} 前 \code{i} 个字符和 \code{s2} 前 \code{j} 个字符，能否交错组成 \code{s3} 前 \code{i+j} 个字符。

转移来自最后一个字符。若 \code{s1[i-1] == s3[i+j-1]}，并且 \code{dp[i-1][j]} 为真，则当前状态可达；若 \code{s2[j-1] == s3[i+j-1]}，并且 \code{dp[i][j-1]} 为真，也可达。两个来源是或关系。开始前必须先判断长度和，否则无论 DP 表如何都不可能匹配。第一行和第一列分别表示只使用一个来源串的情况，遇到第一个不匹配后后面都不能再真。

这题最容易错在把“交错”误解成“保持两串各自相对顺序即可，但不要求用完全部字符”。题目要求刚好用完两个字符串并组成整个 \code{s3}。实验覆盖空串、一个来源为空、两个来源字符大量相同、两条路径都可行、长度不等、最后一个字符只能来自某一侧。空间可以压缩成一维，但压缩后要小心 \code{dp[j]} 代表上一行还是当前行。

\noindent\textbf{LeetCode 115 不同的子序列。}

这题和最长公共子序列很像，但目标不是长度，而是方案数。暴力递归从 \code{s} 中选择或不选择字符，判断能否组成 \code{t}，搜索树会非常大。状态定义为 \code{dp[i][j]}：在 \code{s} 的前 \code{i} 个字符中，有多少个子序列等于 \code{t} 的前 \code{j} 个字符。这里 \code{i} 是候选来源长度，\code{j} 是目标匹配长度，两者角色不能交换。

转移分两种。无论当前 \code{s[i-1]} 是否使用，都可以继承 \code{dp[i-1][j]}，表示不用当前字符；如果 \code{s[i-1] == t[j-1]}，还可以用它匹配目标末尾，增加 \code{dp[i-1][j-1]} 种方案。初始化 \code{dp[i][0] = 1}，因为任何前缀都能用“什么都不选”组成空目标；而 \code{dp[0][j>0] = 0}。

正确性来自最后一个来源字符是否被使用的划分，两类互不重叠且覆盖全部方案。方案数可能很大，C++ 中要根据题目范围考虑 \code{std::int64_t} 或无符号长整型。实验覆盖目标为空、来源为空、来源比目标短、重复字符大量产生多方案、完全无匹配。空间压缩时，\code{j} 必须倒序，否则同一个来源字符会在一轮中被重复使用。

\noindent\textbf{LeetCode 126 单词接龙 II。}

单词接龙 II 不只要求最短长度，还要求输出所有最短路径。暴力 DFS 枚举所有变换路径会陷入指数爆炸，并且容易先走到较长路径。正确思路是先用 BFS 确定最短层级，再用前驱关系回溯所有最短路径。BFS 的层次性保证第一次到达某个单词时是最短距离；但本题必须保留同一层的多个前驱，否则会漏掉另一条同样短的路径。

实现时可以建立通配模式桶，例如 \code{hot} 生成 \code{*ot}、\code{h*t}、\code{ho*}，把共享模式的单词放在同一桶里。这样找邻居时不需要和所有单词逐个比较。BFS 时维护 \code{distance} 和 \code{parents}，当发现一个未访问单词，把距离设为当前距离加一并记录父亲；当发现一个已经在下一层的单词，也要追加当前父亲。不能因为它已经被本层访问就丢掉同层前驱。

BFS 完成后，从终点按 \code{parents} 反向回溯到起点，路径反转后加入答案。边界包括终点不在词典、起点等于终点、多个最短路径、词典中存在大量共享模式、同一层访问标记时机。报告中要说明为什么不能边 BFS 边 DFS 输出路径：那样会混合搜索层级和路径枚举，容易保留非最短路径或漏掉同层父亲。

\noindent\textbf{LeetCode 139 单词拆分。}

单词拆分的暴力递归会枚举每个切分点：如果前缀在字典中，就递归处理后缀。重复后缀会被多次处理，例如很多不同前缀都能走到同一个剩余位置。状态可以定义为 \code{dp[i]}：字符串前 \code{i} 个字符能否被字典拼出。转移枚举最后一个单词的起点 \code{j}，如果 \code{dp[j]} 为真且 \code{s[j:i]} 在字典中，则 \code{dp[i]} 为真。

优化细节有两个。第一，字典放进哈希集合做等值查询，但 \code{substr} 会创建新字符串，热循环里可能增加成本；可以先写清楚版本，再用最大单词长度限制 \code{j} 范围，或用 Trie 避免大量切片。第二，\code{dp[0]} 必须为真，表示空前缀可以被拼出，这是所有第一段单词的起点。如果漏掉这个初始化，任何以开头单词开始的方案都无法转移。

正确性来自最后一个单词。任意合法拆分都有最后一段 \code{s[j:i]}，其前缀 \code{s[0:j]} 必须也可拆；枚举所有 \code{j} 就覆盖全部可能。实验覆盖空串、字典为空、重复单词、长字符串无法拆、多个拆法和最大单词长度剪枝。若题目改成输出所有句子，就不能只保存布尔值，需要在可达性基础上 DFS 枚举路径，并考虑输出规模可能指数级。

\noindent\textbf{LeetCode 188 买卖股票 IV。}

股票 IV 的关键是交易次数。暴力搜索每天买、卖或不操作，会产生大量状态。状态必须包含天数、已完成交易次数和是否持有股票。常见压缩写法维护 \code{buy[t]} 和 \code{sell[t]}：\code{buy[t]} 表示完成至多 \code{t} 次交易后持有股票的最大收益，\code{sell[t]} 表示完成至多 \code{t} 次交易后不持有股票的最大收益。买入不增加完成交易次数，卖出才完成一次交易。

若 \code{k >= n/2}，交易次数限制失去意义，问题退化为无限次交易，可以把所有上涨差价累加，或用两状态 DP。否则每天更新所有 \code{t}。转移时 \code{buy[t] = max(buy[t], sell[t-1] - price)}，\code{sell[t] = max(sell[t], buy[t] + price)}。更新顺序要防止同一天买卖污染状态；一种稳妥方式是使用昨天状态的副本，或者按经过验证的循环顺序更新。

边界包括空价格、只有一天、k 为零、k 很大、连续上涨、连续下跌、价格震荡。正确性来自状态穷尽：每天结束时，要么持有，要么不持有；完成交易次数从零到 k；所有合法操作都在买入、卖出、休息中。面试里要主动说清“交易次数”是约束维度，不是最后统计能补出来的东西。

\noindent\textbf{LeetCode 295 数据流中位数。}

如果每次加入数字后完整排序，插入成本太高。中位数只关心较小一半的最大值和较大一半的最小值，所以用两个堆维护两半。大顶堆保存较小一半，小顶堆保存较大一半；大小差不超过一；并且大顶堆堆顶不大于小顶堆堆顶。两个不变量同时成立时，中位数可以由堆顶直接得到。

插入时可以先根据堆顶把元素放入某一边，再做重平衡；也可以统一先放大顶堆，再把大顶堆堆顶移到小顶堆，最后根据大小移回。无论写法如何，都必须同时维护大小和平序关系。若只维护大小，可能较大元素在左堆；若只维护顺序，大小差过大时堆顶不再代表中位数。

实验覆盖第一个元素、偶数个元素、负数、重复值、升序输入、降序输入、平均值溢出。求偶数中位数时两个 int 相加可能溢出，应先转成 \code{std::int64_t} 或分别除以浮点。若题目扩展为滑动窗口中位数，还要支持删除过期元素，普通堆不支持任意删除，需要懒删除表或有序多重集合。

\noindent\textbf{LeetCode 329 矩阵中的最长递增路径。}

暴力从每个格子 DFS 寻找递增路径，会重复计算同一个后续格子的最长长度。由于路径必须严格递增，状态图是有向无环图：从低值指向高值，不可能回到原值或更低值形成环。可以定义 \code{memo[r][c]} 为从该格子出发的最长递增路径长度。DFS 访问邻居时，只走高度更大的格子，并取邻居结果加一的最大值。

记忆化的正确性来自子问题稳定：从某个格子出发，未来可走路径只由它的位置和矩阵值决定，与之前如何到达它无关。因此计算一次后可以复用。也可以把所有格子按值排序，从大到小或小到大做 DAG DP；本质都是利用严格递增带来的无环性。如果题目允许不下降，就可能出现相等值环，原证明失效。

实验覆盖单格、全相等、严格递增矩阵、严格递减矩阵、多个局部峰值、长蛇形路径。递归实现要注意深度，极端矩阵可能路径很长；工程上可考虑排序 DP 或显式栈。调试时打印 \code{memo} 表，检查每个格子是否只计算一次。

\noindent\textbf{LeetCode 394 字符串解码。}

字符串解码是解析类问题。暴力可以遇到右括号时向前寻找匹配左括号并展开，但嵌套结构会让查找和字符串拼接复杂。更自然的模型是栈保存进入新括号前的上下文：上一层已经构造的字符串、当前重复次数。扫描字符时，数字累积成多位数；遇到左括号，把当前字符串和次数入栈，重置当前字符串；遇到右括号，弹出上一层，把当前片段重复后接回去。

这里的状态必须区分“当前层字符串”和“上一层字符串”。如果只用一个字符串，遇到嵌套时无法知道展开后接回哪里。数字也要累积，例如 \code{12[a]} 不是 \code{1} 和 \code{2} 两个次数。题目通常保证输入合法，但教材实现仍应知道非法输入会在哪里出错：括号不匹配、数字后没有括号、重复次数为空等。

实验覆盖多位数字、嵌套、相邻编码段、普通字符混合、空片段。C++ 中大量字符串重复拼接可能成本高，可以先写清晰版本，再讨论用 \code{reserve} 或输出缓冲优化。这个题和表达式求值同属“上下文栈”：遇到进入嵌套结构时保存旧上下文，退出时恢复并合并。

\noindent\textbf{LeetCode 416 分割等和子集。}

题面问能否把数组分成和相等的两部分。先做代数转换：总和若为奇数，直接不可能；否则问题变成能否选出若干数，使其和为 \code{sum/2}。每个数只能使用一次，所以这是 0-1 背包可达性问题。状态 \code{dp[cap]} 表示是否能用已处理数字凑出容量 \code{cap}。

一维压缩时容量必须倒序遍历。倒序保证当前数字只会基于上一轮状态转移，最多使用一次；如果正序遍历，当前数字刚更新出的状态会在同一轮再次被使用，变成完全背包。这个循环方向不是模板差异，而是物品使用次数的语义。初始化 \code{dp[0] = true}，表示什么都不选可以凑出零。

实验要专门构造循环方向反例，例如只有一个数却能在正序中被重复使用凑出更大容量。边界包括总和为奇数、目标为零、数组中有零、数值很大导致容量表很大。复杂度是伪多项式 \complexity{O(n \cdot target)}，如果数值范围很大，即使元素个数不多也可能慢。

\noindent\textbf{LeetCode 494 目标和。}

每个数前面放正号或负号，暴力枚举所有符号是指数级。设正号集合和为 P，负号集合和为 N，总和为 S，则有 \code{P - N = target}，且 \code{P + N = S}，推出 \code{P = (S + target) / 2}。于是问题变成从数组中选若干数凑出 P 的方案数。这一步转换要先检查 \code{S + target} 是否非负且为偶数，否则无解。

转成背包后，状态 \code{dp[cap]} 表示凑出容量 \code{cap} 的方案数。每个数只能选择一次，所以容量倒序遍历。零元素会自然让方案数翻倍：当 \code{num = 0} 时，\code{dp[cap] += dp[cap]}，表示给零加正号或负号都不改变和，但属于不同符号方案。很多手写特殊分支反而容易错，标准计数转移可以自然处理。

正确性来自代数一一对应。任意符号方案对应一个正号集合，正号集合和必须是 P；任意选出和为 P 的集合，也能给集合内元素正号、集合外元素负号得到目标。实验覆盖目标绝对值大于总和、奇偶不匹配、含多个零、全零、空数组语义。它和分割等和子集很像，但一个求可达性，一个求方案数。

\noindent\textbf{LeetCode 518 零钱兑换 II。}

零钱兑换 II 问组合数，硬币可以无限使用。如果外层遍历金额、内层遍历硬币，会把不同顺序算作不同方案；例如一加二和二加一会被重复统计。组合数要求每组硬币按固定顺序被考虑，因此外层应遍历硬币，内层正序遍历金额。状态 \code{dp[value]} 表示使用已经处理过的硬币种类，凑出金额 \code{value} 的组合数。

初始化 \code{dp[0] = 1}，表示凑出金额零有一种方案：什么都不选。处理某个硬币 \code{coin} 时，金额从 \code{coin} 到 \code{amount} 正序遍历，\code{dp[value] += dp[value - coin]}。正序表示当前硬币可以重复使用；外层硬币顺序固定表示组合不会因为排列不同重复计数。

实验要把循环顺序故意写反，观察组合数如何变成排列数。边界包括金额为零、硬币数组为空、硬币面额大于金额、组合数较大。和 LeetCode 322 的区别必须说清：322 求最少硬币数，518 求组合数量；一个是最值 DP，一个是计数 DP；循环语义和初始化都不同。

\noindent\textbf{LeetCode 621 任务调度器。}

任务调度器可以用堆模拟，也可以用公式。暴力模拟每个时间片，选择当前可执行且剩余次数最多的任务，不够时插入空闲。堆模拟比较直观，但公式能揭示结构：最短时间主要由最高频任务决定。设最高频为 \code{maxFreq}，有 \code{maxCount} 种任务达到最高频。前 \code{maxFreq - 1} 轮需要形成骨架，每轮长度至少为 \code{n + 1}，最后放所有最高频任务的尾部。

公式为 \code{max(tasks.size(), (maxFreq - 1) * (n + 1) + maxCount)}。为什么还要和任务总数取最大？因为其他任务足够多时，它们可以填满所有冷却空位，甚至让整体没有空闲，此时总时间就是任务数。若其他任务不够，骨架长度决定必须插入多少空闲。这个证明比记公式更重要。

实验覆盖冷却为零、只有一种任务、多个最高频任务、其他任务刚好填满空隙、其他任务超过空隙。堆模拟可以作为验证公式的对拍版本。若题目扩展为每个任务有不同执行时间、释放时间或不同冷却间隔，公式不再适用，需要更一般的调度模型。

\noindent\textbf{LeetCode 815 公交路线。}

公交路线题的建模很关键。如果把站点作为 BFS 节点，边表示同一公交路线内任意两站可达，那么一条路线内站点两两连边会非常多。更合适的模型是把“乘坐路线”作为 BFS 层级的一部分：从起点站找到所有可乘路线，乘上一条路线算一次公交；这条路线经过的所有站点都可以到达，并继续找到这些站点可换乘的其他路线。

实现通常建立站点到路线列表的哈希表。BFS 队列里可以放路线，也可以放站点加换乘次数；无论如何，都要分别标记访问过的路线和站点，避免重复展开。同一条路线一旦处理过，就不必再次扫描它的所有站点。目标站如果在当前路线中出现，就可以返回当前乘车次数。

正确性来自 BFS 层次：每一层代表乘坐公交次数增加一，第一次到达目标站就是最少乘车次数。实验覆盖起点等于终点、目标不可达、多条路线共享站点、路线很长、站点编号很大且稀疏。容器选择上，站点编号不一定连续，站点到路线列表适合哈希表；路线访问适合布尔数组，因为路线数量连续。

\noindent\textbf{LeetCode 895 最大频率栈。}

最大频率栈要求弹出出现频率最高的元素；频率相同，弹出最近压入的元素。暴力每次 pop 扫描所有元素找最高频和最近时间，太慢。需要两个维度的索引：value 到频次的哈希表，频次到栈的哈希表，以及当前最大频次 \code{maxFreq}。push 某个值时，先把它的频次加一，再压入对应频次的栈，并更新最大频次。

pop 时，从 \code{maxFreq} 对应的栈弹出栈顶值。由于每次达到某个频次时都按时间压入该频次栈，栈顶自然是该频次下最近的元素。弹出后把该值频次减一；如果最大频次栈为空，\code{maxFreq} 减一。注意，不需要从较低频次栈中删除旧记录，因为那些记录代表该元素历史上达到该频次的时刻；只有当它再次成为当前最大频次栈顶时才会被使用。

正确性来自频次分层和栈内时间顺序。哈希表给出当前真实频次，频次栈保存达到该频次的时间顺序，\code{maxFreq} 指向最高非空层。实验覆盖频次并列、同一元素连续 push、pop 后最大频次下降、多个历史记录存在。这个题能训练“按维度拆索引”：一个结构回答频次，一个结构回答同频最近。

\noindent\textbf{LeetCode 981 基于时间的键值存储。}

题目要求对每个 key 保存多个时间戳版本，并查询不超过给定时间戳的最新值。暴力可以把所有记录放到一个列表里，每次查询扫描；更好的模型是按 key 分组。每个 key 对应一个按时间递增的数组，元素是时间戳和值。查询时先用哈希表定位 key，再在该 key 的时间数组中二分找第一个大于目标时间的位置，答案是它前一个元素。

这个结构利用了题目常见条件：同一个 key 的 set 调用时间戳递增。如果没有这个条件，就需要插入时保持有序，可能使用有序容器或插入后二分位置。二分边界要处理目标时间小于该 key 最早时间，此时返回空字符串；目标时间大于所有记录，返回最后一个值。不要把所有 key 的记录混在一个全局有序数组里，否则查询某个 key 时还要过滤，复杂度和实现都变差。

实验覆盖 key 不存在、时间戳早于所有记录、等于某个记录、落在两个记录之间、晚于所有记录、多个 key 交错 set。C++ 中可以用 \code{unordered_map<string, vector<pair<int,string>>>}，二分时比较 pair 的第一维，或者手写二分避免构造哨兵 pair。这个题的重点是“先按 key 分组，再在组内二分”。

\noindent\textbf{LeetCode 1235 规划兼职工作。}

每个工作有开始时间、结束时间和收益，目标是选择互不重叠工作最大化收益。暴力枚举子集会指数爆炸。排序后可以做 DP：按结束时间排序，\code{dp[i]} 表示考虑前 \code{i} 个工作能获得的最大收益。第 \code{i} 个工作要么不选，收益 \code{dp[i-1]}；要么选，则需要找到最后一个结束时间不大于当前开始时间的工作 \code{p}，收益为 \code{dp[p] + profit[i]}。

二分用于寻找 \code{p}。因为工作按结束时间排序，所有结束时间不大于当前开始时间的工作构成前缀，可以用 \code{upper_bound} 找边界。这里的二分不是查目标值是否存在，而是找兼容工作的数量。若端点语义是前一个工作结束时间等于当前开始时间可兼容，比较条件应使用不大于。

正确性来自最后一个决策：最优解对当前工作只有选和不选两类，选当前工作后，之前能选择的工作只来自兼容前缀，且这部分最优值已经在 DP 中。实验覆盖工作完全不重叠、完全重叠、端点相接、相同结束时间、收益很高但阻塞多个小工作。不要用按收益贪心，局部最高收益可能破坏全局组合。

\noindent\textbf{LeetCode 1438 绝对差不超过限制的最长连续子数组。}

窗口合法性由最大值和最小值决定：窗口内最大值减最小值不超过限制。暴力每个窗口重新找最大最小会很慢；平衡树可以维护动态有序集合，复杂度 \complexity{O(n \log n)}；单调队列可以做到线性。维护两个队列：一个递减队列保存最大值候选，一个递增队列保存最小值候选，队列里存下标以便判断过期。

右端加入新元素时，最大队列弹出所有小于新值的尾部，因为它们既更小又更早，不可能再成为最大值；最小队列弹出所有大于新值的尾部。然后检查队首差值是否超过限制；若超过，就移动左端，并把过期下标从队首弹出。注意，左端可能需要连续移动，直到窗口重新合法。

正确性来自支配关系和窗口单调收缩。队首始终是当前窗口最大和最小；被队尾弹出的旧候选被新候选支配，不会在未来窗口中重新有用。实验覆盖限制为零、全相等、严格递增、严格递减、需要多次收缩、最大最小交替变化。这个题说明滑动窗口状态不一定是和或频次，也可以是动态极值。

\noindent\textbf{LeetCode 1514 概率最大的路径。}

这题是最短路思想的变体。路径概率是边概率乘积，目标是最大乘积。暴力枚举路径不可行；因为边概率在零到一之间，沿路径继续乘只会不增，所以可以用 Dijkstra 类似思想：每次扩展当前概率最大的节点。用大顶堆保存当前到达某节点的最大概率候选，若弹出的概率不是当前记录值，就跳过过期状态。

松弛边时，候选概率是 \code{prob[u] * edgeProb}。如果它大于 \code{prob[v]}，就更新并压入堆。也可以把概率取负对数，把乘积最大转成权重和最小；但直接最大堆更贴近题意。正确性来自单调性：当前堆顶概率最高，未来从其他更低概率路径继续乘边，不可能变成更高概率回头改写它。

实验覆盖起点等于终点、不可达、概率为零、多个同概率路径、浮点比较。C++ 中要使用 \code{double}，并避免用整数处理概率。这个题适合训练“框架迁移”：不是所有 Dijkstra 都叫最短距离，只要代价扩展满足合适单调性，优先队列扩展思想就可能成立。

\noindent\textbf{LeetCode 1675 数组的最小偏移量。}

每个数可以做有限变化：奇数可以乘二一次变成偶数，偶数可以反复除二。直接搜索所有状态组合会爆炸。关键是先把每个数规范化到它能达到的最大偶数形态：奇数乘二，偶数保持。之后每个数只允许不断减小。维护当前所有数中的最大值和最小值，每次尝试降低当前最大偶数，因为只有降低最大值才可能缩小偏移。

用大顶堆保存当前值，同时维护当前最小值。每次堆顶是最大值，先用最大值减最小值更新答案；如果最大值是偶数，可以除二后放回堆，并更新最小值；如果最大值是奇数，它已经不能再降低，过程结束。为什么只处理最大值？偏移由最大和最小决定，降低非最大值不会降低当前最大，反而可能降低最小使偏移变大。

正确性来自状态空间方向。规范化后，每个数的可达序列都是单调下降的；从当前最大值开始下降，枚举了所有可能改善最大边界的步骤。当最大值无法继续下降时，当前及未来都无法缩小最大边界。实验覆盖全奇数、全偶数、已经偏移为零、最大值多次下降、下降后成为新的最小值。注意值乘二可能溢出，要根据输入范围使用足够整数类型。

\noindent\textbf{跨题型迁移和混淆点。}

\noindent\textbf{盛水最多和接雨水不能混用。}

LeetCode 11 盛最多水的容器和 LeetCode 42 接雨水都出现左右指针和高度数组，但问题模型完全不同。盛水最多只选择两条边，面积由短板和宽度决定；每次移动短板，是因为移动长板无法提高短板且宽度必然变小。接雨水则计算每个位置的水量，位置贡献由左右最高边界的较小值决定。前者在候选对之间做排除，后者在每个位置上做边界结算。

如果把盛水最多的证明搬到接雨水，会漏掉中间位置的贡献；如果把接雨水的左右最高思路搬到盛水最多，也会计算不存在的“内部蓄水”。复盘时要写两张状态表：盛水最多保存的是当前左右边界和最大面积，接雨水保存的是左侧最高、右侧最高和当前位置贡献。它们的共同点只是都利用了高度边界，不能因为代码都有双指针就当成同一模板。

\noindent\textbf{最小覆盖子串和找到所有异位词。}

LeetCode 76 最小覆盖子串是可变长度窗口，窗口满足覆盖后要尽量收缩；LeetCode 438 找到字符串中所有字母异位词是固定长度窗口，只需要维护长度等于目标串长度的频次。两者都使用字符计数，但窗口边界语义不同。最小覆盖子串的左端移动由“覆盖条件仍然满足”驱动；异位词窗口的左端移动由“窗口长度固定”驱动。

混用这两题会出现典型错误。把固定窗口写成可变窗口，可能在频次超过需求时过度收缩，漏掉长度刚好的异位词；把可变窗口写成固定窗口，则无法找到更短覆盖。实验时分别打印窗口长度、字符差异和答案更新点。讲题时要先说输出目标：一个要求最短子串，一个要求所有起点；输出目标不同，状态读取位置也不同。

\noindent\textbf{普通 BFS、Dijkstra 和 0-1 BFS。}

看到“最短”不能直接写 BFS。普通 BFS 成立的条件是每条边代价相同，队列层数就是路径长度。Dijkstra 成立的条件是边权非负，每次扩展当前代价最小的未确定状态。0-1 BFS 则针对边权只有零和一的图，用双端队列把零权边放队首、一权边放队尾，从而保持距离顺序。三者都是按某种代价顺序扩展，但容器和证明不同。

网络延迟时间有不同正权边，普通 BFS 会按边数而不是时间扩展，可能错；打开转盘锁每次拨动代价相同，普通 BFS 正合适；带零一代价的状态图若用普通队列，会破坏距离顺序，若用普通 Dijkstra 又能做但常数更高。复盘图题时先写边权范围，再选队列、双端队列或堆。这个判断比背算法名字重要。

\noindent\textbf{并查集和 DFS 连通块。}

岛屿数量、朋友圈、省份数量、冗余连接都和连通性有关，但并查集和 DFS 的职责不同。DFS/BFS 适合从一个起点找完整连通块，能顺便统计面积、边界、路径等信息。并查集适合动态合并和快速回答两个元素是否同组，但它不保存具体路径，也不擅长枚举某个连通块的所有节点，除非额外维护集合成员。

冗余连接按边顺序加入，发现两个端点已经连通时，这条边形成环，因此并查集非常合适。岛屿最大面积需要统计每个连通块大小，BFS/DFS 更直接；当然也可以用并查集维护 size，但实现成本更高。选择结构时先问查询是什么：要遍历块，还是只问同组；要动态加边，还是静态扫描；要路径，还是只要集合代表。

\noindent\textbf{背包三题的状态区别。}

分割等和子集、目标和、零钱兑换 II 都能归到背包，但状态含义不同。分割等和是可达性，\code{dp[cap]} 是布尔值；目标和是方案数，\code{dp[cap]} 是计数；零钱兑换 II 是完全背包组合计数，硬币可无限使用且组合不区分顺序。循环方向和外层顺序都由这些语义决定。

若把分割等和的布尔状态搬到目标和，就只能知道是否存在，无法统计方案；若把零钱兑换 II 的正序容量搬到 0-1 背包，就会重复使用同一个物品；若把组合计数的循环顺序反过来，就会统计排列数。背包题不要背“倒序或正序”，而要先写一句：物品能用几次，答案是可达、最值还是计数，顺序是否有意义。

\noindent\textbf{源码思想迁移到题解表达。}

\noindent\textbf{对象生命周期。}

Linux 源码读多以后，刷题里也会更重视对象生命周期。链表节点是否由题目平台管理，返回的新链表是否复用输入节点，哈希表保存的是值、下标、迭代器还是指针，vector 扩容后旧引用是否还有效，这些都属于生命周期问题。很多面试题不会直接考内存释放，但会通过 LRU、链表反转、合并链表、复制随机链表等题考察你是否理解节点身份。

题解表达中可以主动说明副作用。例如“本题允许修改输入，所以我把访问过的陆地改成水”；“本题不应释放输入链表节点，只需要重连”；“LRU 的哈希表保存链表迭代器，链表节点必须稳定，所以不能用会移动元素的 vector 存节点”。这些话能把算法答案提升到工程答案。

\noindent\textbf{索引组合。}

内核对象常同时进入多个索引：一个对象可能在链表中表示顺序，在树中表示有序键，在哈希表中表示快速查找。刷题中的很多设计题也是索引组合。LRU 是 key 到节点的哈希表加新旧顺序链表；LFU 是 key 到节点、频次到链表、最小频次；时间键值存储是 key 到有序时间数组；账户合并是邮箱到集合编号再按根收集。

讲设计题时，不要只列容器，要说明每个索引回答什么问题。哈希表回答“这个 key 在哪里”，链表回答“谁最旧”，堆回答“当前极值是谁”，有序表回答“前驱后继是谁”，并查集回答“两个对象是否同组”。一个容器若回答不了所有操作，就需要组合；组合后每次更新都要维护所有索引的一致性。

\noindent\textbf{逻辑删除和物理删除。}

RCU、堆懒删除、滑动窗口过期元素、哈希表延迟清理都体现了一个思想：逻辑上不可用不等于物理上马上移除。Dijkstra 堆里的旧状态逻辑上过期，但直到弹出才物理删除；滑动窗口中位数里，过期元素可能埋在堆中，需要延迟清理；RCU 中，旧对象从新索引不可达后，还要等旧读者离开才能释放。

题解里要说清使用前检查。堆顶如果可能过期，就在读取答案前清理；哈希表如果保存历史频次，就要明确什么时候递减或删除；链表节点如果从一个索引摘除，还要确认其他索引是否仍引用它。逻辑删除能降低操作成本，但必须有可靠的清理时机，否则会读到脏状态。

\noindent\textbf{更多实验报告样例。}

\noindent\textbf{样例一：二分答案。}

题目选择 LeetCode 1011 在 D 天内送达包裹。暴力基线是从容量下界开始逐个尝试，检查能否在 D 天内完成。瓶颈是答案空间可能很大，线性试容量太慢。关键观察是容量越大，需要天数越少，可行性单调。检查函数按顺序装包，当前天能放就放，放不下再开新天。这个贪心正确，因为提前开新天只会让剩余天数更紧，不会减少所需天数。

边界测试包括 D 为一、D 等于包裹数、单个包裹等于最大值、容量小于最大包裹、所有包裹相同。实验记录 \code{left}、\code{right}、\code{mid} 和需要天数，检查每次排除是否符合“最小可行容量”语义。C++ 注意总重量可能较大，用足够整数类型。最终复杂度是 \complexity{O(n \log S)}，S 是容量范围。

\noindent\textbf{样例二：单调队列。}

题目选择 LeetCode 239 滑动窗口最大值。暴力基线是每个窗口扫描 K 个元素取最大，复杂度 \complexity{O(nk)}。优化观察是：若新元素比队尾候选更大且下标更晚，那么旧候选在未来任何同时包含二者的窗口中都不可能成为最大值；新元素既更大又更持久。单调队列保存下标，队首是当前窗口最大值。

边界测试包括 K 为一、K 等于数组长度、严格递增、严格递减、重复最大值。调试时打印右端进入、队尾弹出、队首过期和答案读取。证明分两步：队尾弹出不丢答案，因为旧候选被新候选支配；队首过期必须删除，因为它不再属于窗口。C++ 用 \code{deque<int>} 保存下标，而不是保存值，因为需要判断过期。

\noindent\textbf{样例三：树形 DP。}

题目选择 LeetCode 337 打家劫舍 III。暴力递归对每个节点尝试偷或不偷，会重复计算子树。状态应由子树向父节点汇报两个值：偷当前节点的最大收益、不偷当前节点的最大收益。若偷当前节点，孩子不能偷；若不偷当前节点，孩子可以偷也可以不偷，取各自最大值。

边界包括空树、单节点、链式树、左右子树收益差异大。正确性来自树形后序归纳：左右子树已经给出在偷或不偷根时的最优结果，父节点只按相邻不能同偷的约束合并。工程上递归清晰，但极端深树有栈风险；可用显式栈做后序遍历并用哈希表保存每个节点状态。这个实验把 DP 状态和树遍历顺序联系起来。

\noindent\textbf{样例四：区间贪心。}

题目选择 LeetCode 452 用最少数量的箭引爆气球。暴力可以尝试所有射箭位置组合，但选择空间巨大。贪心观察是：按右端点排序后，第一支箭放在最早结束气球的右端，不会比任何其他能射爆这个气球的位置更差，因为它给后续气球留下最大可能覆盖范围。然后跳过所有被这支箭覆盖的气球，继续处理下一个未覆盖气球。

正确性用交换论证。任意最优解中，必须有一支箭射爆最早结束的气球；把这支箭移动到该气球右端，仍能射爆原来被它覆盖且起点不大于该右端的气球，并且不会更早结束。因此存在一个最优解采用贪心第一步。边界包括端点相同、负坐标、区间包含、闭区间端点相接。比较器按右端升序，注意端点可能接近整数边界。

\noindent\textbf{全书复盘索引。}

\noindent\textbf{题型到能力。}

数组训练区间语义和缓存局部性；窗口训练增量维护和单调条件；哈希训练等价类和历史状态；二分训练边界谓词；栈训练最近未完成结构；堆训练动态极值；树训练状态方向；图训练建模和访问标记；DP 训练状态定义和依赖顺序；贪心训练交换证明；并查集训练动态连通；Linux 源码训练对象生命周期和索引组合。

复习时不要按题号孤立记忆，而要问每题训练哪种能力。LeetCode 560 训练前缀代数和哈希频次；LeetCode 862 训练前缀和上的单调队列；LeetCode 1631 训练同一问题的多模型；LeetCode 146 训练索引组合；LeetCode 312 训练区间 DP 的最后一步。把能力写在题号旁边，复习效率会明显高于只刷通过率。

\noindent\textbf{源码到能力。}

Linux 链表训练侵入式节点和稳定身份；hlist 训练桶和删除成本；rbtree 训练动态有序索引；XArray 训练稀疏整数映射；Maple Tree 训练范围查询；RCU 训练逻辑删除和物理释放分离；伙伴系统训练连续性、阶数和合并；SLUB 训练小对象分配成本。它们不是为了替代刷题算法，而是帮助你把数据结构选择讲到工程约束层面。

最终目标是形成统一语言：访问模式决定结构，结构维护不变量，不变量支撑正确性，容器实现带来工程成本，边界实验验证理解。无论面对力扣题、C++ 面试题还是 Linux 源码阅读，都按这条线分析，就不会退回“复制模板”的学习方式。

\noindent\textbf{代码审阅清单。}

\noindent\textbf{状态是否真的必要。}

每写完一道题，先检查状态是否最小且足够。最小是指没有保存未来不会再用的信息；足够是指未来决策所需的信息没有缺失。两数之和只需要历史值到下标，不需要保存所有数对；和为 K 的子数组计数需要历史前缀频次，不是只保存出现与否；最长区间需要最早位置，不是最后位置。状态一旦不匹配输出目标，代码可能在简单样例上通过，却在变体或边界上失败。

审阅时可以问四个问题：状态何时进入，何时更新，何时删除，何时读取答案。窗口题中，右端进入、左端删除和答案读取顺序非常关键；堆题中，堆顶读取前可能需要清理过期元素；DP 题中，状态更新顺序决定读到旧值还是新值。只要这四个问题有一个答不上来，就不应该把题标记为掌握。

\noindent\textbf{复杂度是否按真实维度说明。}

复杂度不能只写一个 \code{n}。图题要区分节点数和边数，字符串 DP 要区分两个字符串长度，背包题要说明容量来自数值大小，答案二分要说明对答案范围取对数，输出所有路径的题要把输出规模算进去。比如单词接龙 II 的 BFS 只是找到最短层级，最终输出所有路径可能本身就很大；N 皇后输出数量也是复杂度的一部分。

审阅复杂度时还要看容器操作。哈希表平均常数不等于最坏常数，\code{substr} 可能复制字符串，\code{priority_queue} 插入和弹出是对数，\code{map} 查询前驱后继也是对数。若代码在循环里构造大量临时字符串，表面 DP 复杂度可能被隐藏成本放大。复杂度说明要和实际代码一致。

\noindent\textbf{边界是否进入主逻辑。}

好的代码不是主体一段、后面一堆补丁分支。边界最好进入初始化和不变量。前缀和把 \code{0 -> 1} 放进哈希表，就自然处理从开头开始的区间；括号栈放 \code{-1} 哨兵，就自然处理第一个有效段；链表使用虚拟头节点，就自然处理删除头节点；二分使用左闭右开，就自然处理插入到末尾。审阅时要优先寻找能消除特殊分支的状态设计。

当然，有些入口边界应当直接返回，例如空数组、起点终点被阻塞、容量为零、目标长度不匹配。这些判断不是补丁，而是题意的不可行条件。区别在于：入口判断处理的是根本无解或无输入；不变量处理的是算法过程中的普通边界。把二者混在一起，会让代码难以解释。

\noindent\textbf{C++ 实现是否隐藏风险。}

C++ 审阅至少看五类问题。第一，整数溢出：面积、路径距离、前缀和、二分平方、概率转换都可能超出 \code{int}。第二，迭代器和引用失效：\code{vector} 扩容、\code{erase}、\code{unordered_map} rehash 都可能让保存的引用失效。第三，容器副作用：\code{operator[]} 会插入默认值。第四，比较器：排序和堆比较器必须满足严格弱序。第五，复制成本：循环中不要无意复制大 \code{vector}、\code{string} 或复杂节点。审阅必须具体。

设计类时还要看成员职责。并查集的 \code{find} 是否路径压缩，\code{unite} 是否按大小合并；LRU 的移动节点是否只改链表而不丢哈希；Trie 的析构和节点所有权是否清楚；图算法的邻接表是否按输入规模预分配。刷题代码不一定要写成工程库，但高质量教材必须告诉读者这些风险在哪里。

\noindent\textbf{十周训练安排。}

\noindent\textbf{第一、二周：线性结构。}

第一周集中数组、字符串、双指针和窗口。每天选择两道基础题和一道变体题，必须写暴力版和优化版。重点不是数量，而是把区间语义写清楚：左闭右开还是左闭右闭，慢指针左侧保存什么，窗口何时合法，答案何时更新。实验必须打印指针移动过程。推荐题包括两数之和、无重复字符最长子串、盛水最多、三数之和、最小覆盖子串、长度最小子数组、找到所有异位词。

第二周集中链表、栈、队列和单调结构。链表题每一步都画前驱、当前、后继；栈题说明栈顶代表什么；单调栈和单调队列说明支配关系。推荐题包括反转链表、K 个一组翻转链表、有效括号、最长有效括号、每日温度、柱状图最大矩形、滑动窗口最大值。周末写一个 LRU，把哈希表和双向链表的职责讲清楚。

\noindent\textbf{第三、四周：查找和有序结构。}

第三周集中哈希、前缀和和二分。哈希题必须写 key 和 value 的语义，前缀题必须写代数等式，二分题必须写谓词和区间不变量。推荐题包括字母异位词分组、和为 K 的子数组、连续子数组和、连续数组、搜索插入位置、查找首尾位置、爱吃香蕉、船运包裹、分割数组最大值。

第四周集中排序、区间、堆和有序容器。排序题说明排序换来了什么关系，区间题说明端点闭开，堆题说明堆顶是否可能过期，有序容器题说明前驱后继。推荐题包括合并区间、插入区间、无重叠区间、用箭引爆气球、前 K 高频、数据流中位数、最接近原点 K 个点、任务调度器。周末写 map 动态区间实验，对照 Maple Tree 思维。

\noindent\textbf{第五、六周：树和图。}

第五周集中二叉树、BST、Trie。每道树题先判断状态方向：自顶向下、后序汇总还是层序传播。递归写法要能说出迭代替代方案。推荐题包括层序遍历、验证 BST、最大深度、最近公共祖先、二叉树直径、打家劫舍 III、单词搜索 II、实现 Trie。实验打印显式栈和队列状态。

第六周集中图、拓扑、并查集和最短路。每道图题先写节点、边、边权和状态维度。推荐题包括岛屿数量、岛屿最大面积、腐烂橘子、打开转盘锁、课程表、省份数量、冗余连接、网络延迟时间、最小体力路径、概率最大路径。周末比较 BFS、Dijkstra、二分 BFS 和并查集在同一题上的建模差异。

\noindent\textbf{第七、八周：动态规划。}

第七周集中线性 DP、二维 DP 和字符串 DP。每题先写暴力递归，再圈出重复状态。推荐题包括爬楼梯、打家劫舍、最大子数组和、不同路径、最小路径和、最长公共子序列、编辑距离、交错字符串、不同子序列。每题至少打印一个小 DP 表，解释每个格子来源。

第八周集中背包、区间 DP、状态机和状态压缩。推荐题包括零钱兑换、零钱兑换 II、分割等和子集、目标和、买卖股票含冷冻期、买卖股票 IV、戳气球、最长递增子序列。重点检查循环方向、初始化、空间压缩旧值来源和伪多项式复杂度。周末把三道背包题放在一张表里比较状态含义。

\noindent\textbf{第九、十周：综合和源码。}

第九周做综合设计题和高频难题。推荐题包括 LFU 缓存、最大频率栈、时间键值存储、公交路线、最短子数组至少为 K、数组最小偏移量、规划兼职工作。每题都写索引组合图，说明每个结构回答什么查询，更新时如何保持一致。

第十周读源码和复盘。按链表、hlist、rbtree、XArray、Maple Tree、RCU、伙伴系统、SLUB 的顺序，每天只读一个小问题，写一页报告。最后选十道自己最容易混淆的题，重新讲一遍暴力、瓶颈、状态、不变量、边界和 C++ 风险。十周结束后，不要求记住所有题号，但要能看到题面就判断访问模式、状态维度和可疑边界。

\noindent\textbf{质量底线。}

\noindent\textbf{不能用模板替代理解。}

任何题解如果只有“这题用哈希”“这题用 DP”“这题用双指针”，都不合格。必须补上为什么：哈希保存什么等价类，DP 状态包含哪些未来信息，双指针排除了哪一批候选。模板只能作为复习索引，不能作为解释本身。读者应该在每道题里看到暴力版本如何演化，而不是看到一个突然出现的最终答案。

教材写作也一样。章节之间要有衔接：数组的区间语义过渡到窗口，窗口的增量维护过渡到单调队列，前缀和的代数关系过渡到哈希历史状态，树的后序汇总过渡到树形 DP，图的层次扩展过渡到最短路。只有这样，书才像教材，而不是题解堆。

\noindent\textbf{不能用字数掩盖空洞。}

字数目标只是底线，不是质量本身。有效内容应当增加背景、推导、证明、边界、实验或源码对照；无效内容只是换词重复同一句话。审阅时要删掉不能回答新问题的段落。一个段落如果不能说明“为什么这样做”“错在哪里”“边界如何处理”“和另一个结构有什么区别”，就应该重写。

本书后续继续扩展时，应优先补三类内容：第一，更多经典题的完整推导；第二，真实 C++ 实验和性能对拍；第三，Linux 源码阅读报告。它们能自然增加深度，不需要硬凑。最终读者拿到的应是一套可执行的学习系统：能刷题、能讲题、能写代码、能读源码。

\noindent\textbf{口述和代码审查样例。}

\noindent\textbf{样例一：讲清前缀和。}

以前缀和题为例，合格口述不能停在“我用前缀和”。应当这样讲：我先把任意区间和写成两个前缀的差，因此当前右端固定时，问题变成寻找某个历史前缀。若题目问数量，历史表保存频次；若题目问最长，历史表保存最早位置；若题目问是否存在，保存布尔即可。这个区分直接决定代码里的 value 类型。

审查代码时，第一看初始化是否包含前缀零；第二看先查询还是先更新；第三看是否使用足够整数类型；第四看负数是否影响算法。和为 K 的子数组不需要数组全正，因为它只依赖代数等式；长度最小且至少为 K 的问题若含负数，普通窗口失效，要换成前缀和加单调队列。能说清这两个题的差异，就说明前缀和不是死记公式。

\noindent\textbf{样例二：讲清二分。}

二分题的口述要先定义谓词。比如船运包裹，谓词是“容量为 cap 时，能否在 D 天内按顺序运完”。这个谓词随 cap 增大从假变真，因此要找第一个真。区间语义可以固定为左闭右开：答案始终在 \code{[left,right)} 内，若 \code{mid} 可行，就把右边界收到 \code{mid}；否则把左边界移到 \code{mid+1}。

代码审查时，不要只看循环是否会停。要看下界是否为最大包裹，上界是否为总重量，检查函数是否在容量小于某个包裹时正确失败，求和是否溢出。很多二分错题的根源不在二分循环，而在谓词写错。二分只是加速寻找边界，不能修复错误的可行性判断。

\noindent\textbf{样例三：讲清单调栈。}

单调栈题的口述要说明“栈里保存的是还没有被解决的候选”。每日温度中，栈里是还没找到更高温的日期；柱状图中，栈里是还没确定右侧第一个更矮位置的柱子；下一个更大元素中，栈里是等待未来更大值的元素。新元素到来时，如果它能解决栈顶，就不断弹出并结算。

代码审查时，第一看栈里存值还是下标。只要需要距离、宽度或过期判断，就必须存下标。第二看相等元素如何处理。第三看是否需要哨兵。第四看弹出时结算对象是谁。柱状图最大矩形里，当前更矮柱不是答案高度，它只是右边界；被弹出的柱子才是本次结算的最矮高度。

\noindent\textbf{样例四：讲清图建模。}

图题口述第一句必须定义节点和边。岛屿数量的节点是陆地格子，边是四方向相邻；课程表的节点是课程，边是先修关系；公交路线题更适合把路线或“乘坐路线”作为 BFS 层次，而不宜把同一路线内所有站点两两连边；最小体力路径的边权是高度差，路径代价是最大边权，不是和。

代码审查时，第一看访问标记时机。BFS 通常入队时标记，避免重复入队。第二看状态维度。位置相同但钥匙集合不同、剩余障碍次数不同，都不能合并。第三看边权。等权用 BFS，非负权用 Dijkstra，零一权用双端队列。第四看邻居生成成本。单词接龙如果每次扫描所有单词找邻居，复杂度会很差，应考虑通配桶。

\noindent\textbf{样例五：讲清 DP。}

DP 口述要从暴力搜索树开始。先说每个搜索节点需要哪些信息，再说哪些节点会重复出现，最后给状态命名。编辑距离的状态是两个前缀长度；背包的状态是物品阶段和容量；股票的状态是天数、持有状态和交易次数；树形 DP 的状态是子树向父节点返回的信息。状态不是为了填表而填表，而是为了让未来决策只依赖少量参数。

代码审查时，第一看初始化是否对应空状态；第二看转移是否覆盖所有最后一步；第三看遍历顺序是否满足依赖；第四看空间压缩是否读旧值。0-1 背包倒序容量，完全背包正序容量，LCS 一维压缩保存左上角旧值，股票状态机保存昨天状态。只要空间压缩解释不清，就先退回完整二维或多状态表。

\noindent\textbf{源码阅读样例报告。}

\noindent\textbf{报告一：链表删除。}

问题：删除一个已知链表节点需要维护什么不变量。对象：业务结构体中嵌入一个链表节点。索引：链表按某种顺序组织对象。插入时要让前驱、后继和新节点三者互相指向；删除时要让前驱直接指向后继，并让后继回指前驱。若遍历中删除当前节点，必须先保存下一个节点，否则删除后当前节点的链接可能被重置。

刷题类比：LRU 缓存移动节点到头部，K 个一组翻转链表接回前后段。工程差异：内核对象可能同时挂在多条链上，删除某条链的节点不代表对象释放。并发情况下，还要考虑锁或 RCU。实验结论：链表操作的难点不是复杂度，而是节点身份、链接顺序和生命周期。

\noindent\textbf{报告二：红黑树插入。}

问题：动态有序索引如何插入新对象。对象：业务结构体中嵌入树节点。索引：按某个 key 维护二叉搜索树顺序。插入前，调用者沿树比较 key，找到空孩子位置；链接新节点后，库函数通过旋转和变色修复平衡。旋转不改变中序顺序，所以查找语义保持不变。

刷题类比：日程安排用有序表找前驱后继，数据流排名维护动态有序集合。工程差异：C++ \code{map} 隐藏节点分配和比较器，Linux rbtree 让调用者负责比较和对象嵌入。实验结论：有序树解决的是顺序查询，不是比哈希表更高级；只要题目需要前驱后继，哈希表就不够。

\noindent\textbf{报告三：RCU 删除。}

问题：读多写少结构中，删除对象为什么不能立刻释放。对象：读者可能正在访问的共享节点。索引：链表、树或哈希表让新读者找到对象。删除第一步是从索引摘除，让后续读者不可达；第二步等待旧读者离开读侧临界区；第三步释放内存。若摘除后马上释放，旧读者可能访问悬空指针。

刷题类比：Dijkstra 堆中的旧状态、滑动窗口堆中的过期元素。相同点是逻辑不可用和物理清理分离；不同点是 RCU 为并发安全服务，堆懒删除多为容器能力限制服务。实验结论：删除语义要拆开看，不能把“从答案集合移除”和“销毁对象”混成一步。

\subsection{单点深挖和高频设计}

\noindent\textbf{单点深挖：把模板拆成问题。}

这一节专门回应一个写作底线：单点问题要单点分析，不能把同一套话粘到每道题下面。真正的教材应该让读者看到每个问题自己的结构。接雨水不是“又一个双指针”，它的难点是何时可以结算某个位置的水；柱状图不是“又一个单调栈”，它的难点是被弹出的柱子在那一刻同时知道左右边界；LRU 不是“哈希加链表”，它的难点是节点身份、两个索引同步和对象生命周期。下面每个小节只抓一个具体难点，把暴力、优化、反例、实验和工程迁移连成一条线。

\noindent\textbf{LeetCode 42 接雨水：先问结算对象。}

接雨水最常见的错误讲法是直接背“左右指针，谁小移谁”。这句话不解释为什么能算水，也不解释什么时候不能算。先从暴力出发：对位置 i，水位等于左侧最高柱和右侧最高柱的较小值，再减去当前位置高度。暴力每个位置都向两边扫描，所以慢在反复寻找左右边界。前后缀数组的优化很直观：左边界数组保存每个位置左侧最高，右边界数组保存每个位置右侧最高，最后逐点结算。这个版本空间高一些，但最容易证明，因为每个位置结算时两个边界都已经明确。

双指针版本必须继承这个结算条件。设左侧已经看到的最高柱是 left max，右侧已经看到的最高柱是 right max。如果 left max 不大于 right max，那么左指针位置的右侧至少存在一个不低于 left max 的边界，所以它的水位上限已经由 left max 决定，此时可以结算左侧并右移。反过来，如果 right max 更小，就结算右侧。注意，结算的是较低已知边界那一侧的位置，不是简单因为某个指针当前高度小就移动。代码里可以用当前高度更新边界，再根据边界大小结算，也可以先比较再更新，但不变量必须保持一致。

这道题的实验要同时跑三版：暴力版、前后缀版、双指针版。输入选择单调递增、单调递减、平台高度、多个凹槽、长度不足三、两端高而中间低。打印每一轮 left、right、left max、right max 和本轮新增水量。若某个实现提前结算了右侧边界未知的位置，平台高度和多凹槽会立刻暴露问题。把这组实验做完，读者才能明白双指针不是口诀，而是把“每个位置需要左右最高”压缩成两个移动边界。

单调栈解法是另一种结算对象。栈中保存还没有找到右边界的柱子，且高度保持单调。当遇到更高的当前柱时，栈顶作为槽底被弹出，当前柱是右边界，弹出后的新栈顶是左边界。此时槽底水量才确定。这个版本特别适合解释“弹出时机”：不是当前柱接水，而是被弹出的低柱终于知道了左右边界。把双指针和单调栈放在一起比较，能训练一个重要能力：同一题可以有不同状态，但每个状态都必须回答“现在能结算谁”。

\noindent\textbf{LeetCode 76 最小覆盖子串：覆盖不是匹配总数。}

最小覆盖子串的暴力做法是枚举所有子串，再检查是否覆盖目标字符需求。这个暴力版清楚地暴露两个重复：窗口候选太多，检查需求时又反复统计。滑动窗口能成立，是因为右端扩展只会增加窗口里的字符，覆盖条件更容易满足；当窗口已经覆盖目标时，继续扩展只会更长，所以应当收缩左端尝试变短。这个推理依赖“覆盖”是一个可增量维护的条件。

很多错误实现把覆盖理解成匹配字符总数。目标串如果是 AABC，只看到一个 A 不能算完成 A 的需求；窗口里多出来的无关字符也不能增加完成度。更稳的状态是 need 保存目标每个字符所需次数，window 保存当前窗口次数，formed 表示有多少种字符已经达到需求。一个字符从少于需求变成恰好达到需求时，formed 加一；从恰好达到需求被左端移走变成不足时，formed 减一。formed 等于 need 中字符种类数，才表示窗口覆盖目标。这里的“种类达标”比“总匹配数”更不容易被重复字符误导。

证明时按固定右端来看。当右端扩展到覆盖目标，内层循环会把左端一直右移到再移就不覆盖为止。在这个过程中，每个合法窗口都被检查，最短的那个被记录。对所有右端重复这个过程，全局最短一定被看见。被移过的左端不会再产生更短答案，因为未来右端只会更大，窗口只会更长。这个排除逻辑必须讲清楚，否则读者只能记住一段滑动窗口代码。

实验输入要包含目标字符重复、目标字符缺失、答案在开头、答案在结尾、大小写混合、源串短于目标串。调试时打印 left、right、当前字符、formed、need size 和当前最优区间。若用总匹配数写法，目标含重复字符的输入会暴露错误；若左端收缩后忘记更新 formed，答案在开头或结尾的用例会出错。C++ 里字符集固定时用数组即可；若题目不保证字符集，再用哈希表。容器选择只是实现细节，核心仍是“覆盖状态”到底表示什么。

\noindent\textbf{LeetCode 862 最短子数组至少为 K：负数为什么击穿窗口。}

很多人看到“连续子数组”和“至少为 K”就想用滑动窗口。若数组全为正数，右端扩展会让窗口和变大，左端收缩会让窗口和变小，所以窗口条件具有单调性。但一旦允许负数，右端加入一个负数可能让和变小，左端移走一个负数可能让和变大，左端不再能安全排除。用一个小反例就能说明问题：窗口当前不满足 K，不代表继续扩展一定更接近；窗口当前满足 K，移走左端也不一定破坏或改善。普通窗口失去依据。

这题正确的切入点是前缀和。区间和等于 prefix[j] 减 prefix[i]，目标是找到 j 固定时最靠右、同时 prefix[i] 不大于 prefix[j] 减 K 的历史下标 i。为了让长度最短，历史下标越大越好；为了让未来更容易满足条件，前缀和越小越好。单调队列保存候选历史下标，队列里的前缀和递增。扫描当前前缀时，若当前前缀减队首前缀至少为 K，就可以更新答案，并弹出队首，因为对未来来说，更靠后的当前 j 已经让这个队首不可能给出更短长度。随后，若队尾前缀和大于等于当前前缀，则队尾被当前下标支配：当前下标更靠后，前缀和还不大，未来任何右端选择当前下标都不差。

这道题的重点是两个弹出条件含义完全不同。队首弹出是“已经形成一个答案，为了找更短长度继续丢旧左端”；队尾弹出是“旧候选被新候选支配，以后没必要保留”。如果只背单调队列，很容易把两个条件混成一个。实验要准备含负数、全正、答案不存在、答案为一、前缀和反复下降、K 很大这些输入。打印 prefix、队列下标、队列前缀值和答案。看到队尾被支配的过程，读者才会理解单调队列不是排序结构，而是保留一组仍有可能成为最优左端的历史状态。

这题也适合作为“模板失效”的案例。LeetCode 209 长度最小的子数组可以用正数窗口，LeetCode 560 和为 K 的子数组可以用前缀和加哈希计数，LeetCode 862 要用前缀和加单调队列。三个题都在说连续区间，但保存的历史状态不同：正数窗口保存当前和，计数题保存历史前缀频次，最短且含负数保存单调的历史前缀下标。把它们放在一起，读者才能学会根据条件换结构。

\noindent\textbf{LeetCode 84 柱状图：弹栈时到底算谁。}

柱状图最大矩形的暴力模型很清楚：把每根柱子当作矩形最矮高度，向左右扩展到第一个更矮柱子之间，面积就是高度乘宽度。暴力慢在每根柱子都重复找左右第一个更矮位置。单调栈的目的就是在一次扫描中确定这些边界。栈中保存高度递增的下标，表示这些柱子还没遇到右侧第一个更矮位置。

当当前柱子高度小于栈顶柱子时，栈顶柱子的右边界已经确定为当前下标。弹出栈顶后，新的栈顶就是它左侧第一个更矮柱子。于是被弹出的柱子作为最矮高度时，最大宽度就是右边界减左边界再减一。如果弹出后栈为空，说明左侧没有更矮柱，宽度从零延伸到当前下标之前。这里最容易错的是把当前柱子当成答案高度。当前柱子只是触发结算的右边界；真正被结算的是被弹出的柱子。

相等高度也要有一致策略。可以遇到小于栈顶才弹，也可以遇到小于等于时弹，只要宽度语义和重复高度处理一致。教材阶段建议先用末尾哨兵零，让所有柱子最终都被结算，避免循环结束后再写一段补丁逻辑。实验输入必须包括严格递增、严格递减、全相等、单柱、空数组和中间凹槽。严格递增会暴露末尾未结算，严格递减会暴露宽度计算，全相等会暴露相等高度策略。

这道题还连接 LeetCode 85 最大矩形。二维矩阵中每一行都可以看成柱状图底部，连续一的高度向下累积，然后对每一行运行柱状图算法。这里的迁移不是复制代码，而是迁移“某个高度作为最矮高度时找左右边界”的模型。读者若能从柱状图自然推到最大矩形，就说明已经掌握了单调栈的结算对象，而不是只记住一段弹栈代码。

\noindent\textbf{LeetCode 15 和 18：排序、去重和输出规模。}

三数之和最容易被讲成“排序加双指针”，但排序为什么有用，去重为什么必须分层，往往没有讲透。暴力三层循环覆盖所有下标三元组，慢在组合数量太多，也慢在重复值会生成大量相同答案。排序后的第一层收益是制造单调性：固定第一个数后，剩下两个数的和随左指针右移而增大，随右指针左移而减小。这样可以安全排除一批配对，而不是逐个枚举。

第二层收益是去重变得局部可见。固定值如果和前一个固定值相同，跳过当前固定值，因为它会生成同一组数值答案。找到一组三元组后，左指针和右指针也要跳过相同值，避免相同数值组合重复进入结果。注意，去重不是最后把答案塞进集合。集合过滤能让结果看起来正确，但它掩盖了算法结构，也增加了额外成本。高质量题解应该在生成阶段就避免重复答案。

四数之和是三数之和的压力测试。外层多固定一个数，目标和可能超出 int，答案数量也可能很大。排序后可以用最小可能和与最大可能和提前剪枝：若当前固定前缀加上后面最小几个数已经大于目标，后续更大固定值也不必继续；若加上最大几个数仍小于目标，则当前固定值太小，可以跳过。这个剪枝必须建立在排序上，也必须注意溢出。C++ 中目标和、中间和最好用更宽整数。

实验要包含全零、大量重复、无答案、正负混合、目标极大、答案很多。还要专门测输出规模：当答案本身很多时，复杂度不能只写计算部分，结果数组也占空间和时间。面试追问 k 数之和时，不要直接写递归模板。先说明递归每层固定一个数，最后降到双指针；再说明每层去重和剪枝如何保持不重复、不漏解。这样讲，三数、四数和 k 数才是一条连续的教材线。

\noindent\textbf{答案二分：真正要审查的是谓词。}

二分答案题表面上在写循环，实际难点在检查函数。以 LeetCode 410 分割数组的最大值为例，猜一个最大段和上限 limit，检查能否把数组按原顺序分成不超过 k 段，使每段和不超过 limit。检查函数的贪心是尽量把当前数字放进当前段，放不下才开新段。为什么这样正确？因为顺序固定，提前开新段只会让剩余段数更紧，不会减少所需段数。这个证明比外层二分更重要。

LeetCode 875 爱吃香蕉也是一样。速度 speed 越大，吃完所需时间越少，因此可行性从假到真。检查函数里每堆耗时要向上取整，不能用普通除法。边界也由谓词决定：速度下界是一，上界可以取最大堆；船运包裹的容量下界应是最大包裹，上界是总重量；分割数组的 limit 下界是最大元素，上界是总和。若下界设错，检查函数可能处理本不该出现的非法状态。

二分实验要把外层循环和检查函数分开测。先写一个线性枚举答案的慢版本，再写二分版本，两者对拍。对每个 mid 打印 left、right、谓词值和检查函数返回的段数或小时数。若二分结果错，先怀疑谓词，不要立刻改加一减一。很多错误都是“谓词并不单调”或者“检查函数没有表达题意”，外层循环再标准也无济于事。

这类题还要说明复杂度维度。它通常不是 \complexity{O(n\log n)}，而是 \complexity{O(n\log M)}，其中 M 是答案范围。背包、金额、容量、速度、最大段和这类数值进入复杂度时，要主动说明它和输入长度不是同一个维度。面试中能讲清这个点，会明显区别于只背模板的人。

\noindent\textbf{LeetCode 146 LRU：两个索引和一个节点身份。}

LRU 缓存的暴力实现可以用数组保存访问顺序。每次 get 或 put 时线性查找 key，再把元素移动到头部；容量满时删除尾部。这个版本正确但慢，瓶颈在定位 key 和移动节点。哈希表能常数定位 key，但不能表达最近使用顺序；双向链表能常数删除已知节点并移动到头部，但不能按 key 查找。组合结构的必要性就在这里：哈希表负责定位，链表负责顺序，二者共享同一个节点身份。

节点里必须保存 key 和 value。很多人疑惑哈希表已经有 key，链表节点为什么还要保存 key。原因在淘汰尾部时，链表先告诉我们最久未使用的节点，但要从哈希表中删除对应条目，必须知道它的 key。若节点只保存 value，就无法常数时间同步删除哈希索引。更新已有 key 时也不能新建重复节点，否则哈希表和链表会出现两个身份不一致的对象。正确做法是更新节点值，并把同一个节点移动到头部。

这题可以直接对照 Linux 侵入式链表。刷题版常用 \code{std::list} 让容器拥有节点，哈希表保存迭代器。内核常把 \code{list_head} 嵌进业务对象，链表只组织对象，不拥有对象。共同点是节点身份必须稳定，删除必须同步多个索引；区别是内核还要处理引用计数、锁、RCU 和分配上下文。把这种差异讲出来，LRU 就不再是六个字，而是一个小型索引系统。

实验要跑操作序列而不是只测单次 get。容量为一、重复 put、get 不存在、更新已有值、get 后再 put、连续淘汰都要覆盖。每步打印链表从头到尾的 key 顺序和哈希表大小，检查两个不变量：哈希表每个 key 都能在链表中找到唯一节点，链表尾部就是最久未使用节点。若某个操作后链表有重复 key 或哈希表指向失效迭代器，说明节点身份没有维护好。

\noindent\textbf{并查集：代表元不是路径。}

并查集适合回答“两个元素是否在同一集合”以及“把两个集合合并”。它不保存两点之间的具体路径，也不适合回答最短距离。LeetCode 684 冗余连接中，按顺序加入无向边；若一条边的两个端点已经同属一个集合，再加入它就会形成环。这道题不需要知道环上所有边，也不需要给出路径，所以并查集正好匹配。暴力每加一条边都 DFS 检查连通性，会重复扫描已有结构。

并查集的正确性来自等价关系。每个集合有代表元，find 返回代表；unite 只把两个代表连起来。路径压缩改变的是树形实现，不改变集合语义；按大小合并降低树高，也不改变集合语义。初学者容易把 parent 数组看成图的真实父边，这是错误的。压缩后 parent 会直接指向根，原图路径信息已经不存在。如果题目要输出路径，不能只靠普通并查集。

账户合并是另一个代表性例子。邮箱是连接账户的证据，姓名通常不能作为合并依据。可以给每个邮箱编号，把同一账户中的邮箱合并；最后按根收集邮箱并排序。这里并查集解决的是动态连通，哈希表解决的是字符串到编号的映射，排序解决的是输出格式。三个结构各司其职。若题目变成等式除法或带奇偶约束，并查集就要扩展权值或颜色；普通代表元不够。

实验要打印合并前后每个元素的根、集合大小和 parent 树高度。重复合并、已经连通的边、孤立点、字符串编号映射都要测。还可以故意去掉按大小合并或路径压缩，观察 find 路径变长。这个实验的目的不是追求近乎常数的理论证明，而是让读者看见代表元如何把“许多连接关系”压缩成一个可查询的集合状态。

\noindent\textbf{Dijkstra：旧堆状态不是错误。}

带正权图不能用普通 BFS，因为 BFS 按边数扩展，不按总代价扩展。LeetCode 743 网络延迟时间中，从起点到各点的最短时间取决于边权和。Dijkstra 的核心状态是当前已知最短距离，优先队列每次弹出距离最小的候选。若边权非负，弹出的未过期状态已经不可能被其他路径再改小，因此可以扩展它的出边。

C++ 的 \code{priority_queue} 没有 decrease-key。发现更短距离时，常见写法不是修改堆中旧元素，而是把新距离再压入堆。这样堆里会同时存在同一个节点的旧状态和新状态。弹出时比较堆顶距离和 distance 数组，若不一致，说明它已经过期，跳过即可。旧堆状态不是错误，只要使用前检查，它只是容器接口限制下的工程实现。

这点可以和 RCU 或懒删除做概念对照，但不能混为一谈。Dijkstra 的旧状态是算法实现中的过期候选，跳过即可；RCU 的旧对象是并发读者可能仍在访问的内存，必须等宽限期后释放。共同点是逻辑状态和物理存在分离；不同点是安全目标完全不同。教材要把相似性和差异都说清楚，避免把源码概念泛化成空话。

实验要构造重边、零权边、不可达节点、多条等长路径和明显会产生旧堆状态的图。打印每次弹出的节点、堆中距离、当前 distance、是否过期。再用普通 BFS 跑同一张带权图，观察结果差异。最后让读者回答：如果边权为负，为什么这个证明失效？这个问题能把 Dijkstra 的条件边界讲清楚。

\noindent\textbf{DP：从暴力搜索树命名状态。}

动态规划最怕直接给公式。以编辑距离为例，暴力思路是从两个字符串末尾开始，如果末尾字符相同，就处理剩余前缀；如果不同，就尝试删除、插入、替换。这个递归会反复处理同一对前缀长度。于是状态自然命名为 dp[i][j]，表示第一个字符串前 i 个字符变成第二个字符串前 j 个字符的最少操作。公式不是凭空出现，而是暴力选择的缓存。

背包也是同样逻辑。暴力枚举每个物品选或不选，未来只依赖“处理到第几个物品”和“剩余容量”。0-1 背包每个物品只能用一次，所以一维压缩时容量倒序，防止当前物品在同一轮被重复使用；完全背包允许重复使用当前物品，所以容量正序。若把循环方向背成模板，不构造反例，就很容易在分割等和子集或零钱兑换之间写错。

最长递增子序列展示了另一类优化。二次 DP 的状态是必须以 i 结尾的最长长度，枚举所有前驱。二分优化维护 tails，表示某个长度的递增子序列能拥有的最小结尾。tails 不是答案序列，而是未来扩展潜力表。这个区别必须说清楚，否则读者会误以为 tails 中的值就是实际选择路径。若要还原路径，还要记录前驱和每个长度对应的位置。

DP 实验不要一上来压缩空间。先打印完整表，让读者解释每个格子来自哪里。编辑距离表能看到插入、删除、替换；LCS 表能看到不连续子序列；背包表能看到每个物品阶段如何影响容量；股票状态机能看到持有、卖出、休息的约束。只有完整表能讲清，再做滚动数组。空间压缩是优化，不是理解的起点。

\noindent\textbf{Trie 和单词搜索：前缀剪枝不等于路径合法。}

单词搜索 II 的暴力做法是对每个单词在棋盘上单独 DFS。慢点在于不同单词共享大量前缀，却被重复搜索。Trie 把所有单词的前缀合并，棋盘 DFS 时沿 Trie 节点前进；如果当前路径在 Trie 中没有对应前缀，就立即剪枝。这个优化节省的是“重复检查前缀”的工作，而不是改变网格路径规则。

路径合法性仍然由访问标记控制。同一条路径中，一个格子不能重复使用；进入格子要标记，递归返回要恢复。Trie 节点上存在单词结束标记，只表示当前路径构成一个词；它不允许你忽略网格访问约束。找到一个词后，可以清空该节点的单词标记避免重复加入答案，也可以在工程上做叶子删除减少搜索，但这些都不能破坏其他单词共享的前缀。

实验要准备重复字母棋盘、单字符单词、一个单词多条路径、多个单词共享前缀、路径必须回退的用例。打印当前坐标、路径字符串、Trie 深度和访问矩阵。若忘记恢复访问标记，其他分支会漏答案；若没有访问标记，同一格子会被重复使用；若找到词后错误删除共享节点，另一个共享前缀的词会丢失。这个单点能把回溯、Trie 和状态恢复三者联系起来。

C++ 实现时要考虑节点所有权。教学代码可以用数组下标池或智能指针；若用裸指针，必须说明释放责任。字符集固定为小写字母时，子节点数组更快；字符集稀疏时，哈希表更省空间。不要在 DFS 热路径里反复构造大字符串，可以用路径数组或直接在 Trie 节点保存完整单词。容器选择仍然服务于访问模式。

\noindent\textbf{动态区间和 Maple Tree 思维。}

离线区间题和动态区间题的难点不同。LeetCode 56 合并区间只需排序后线性扫描；所有区间一次性给出，排序把可能相交的区间放到相邻位置。日程安排、内存区间表或虚拟地址区域管理则是动态问题：区间会插入、删除、查询，不能每次都重新排序全部数据。此时需要有序结构维护相邻关系。

用 \code{std::map} 写简化动态区间表，key 是区间起点，value 是终点和属性。点查询时，找起点不大于目标地址的最后一个区间，再判断终点是否覆盖。插入新区间时，只需检查前驱和后继是否重叠；如果它们都不冲突，更远区间也不可能冲突。删除一段时更复杂：可能删除整个区间，可能截掉左侧或右侧，可能从中间挖掉一段并分裂成两个区间。

Maple Tree 的源码背景比这个实验复杂得多，但问题动机相通：大地址空间、稀疏区间、频繁查询、动态更新、读路径性能和并发安全。刷题阶段不需要复刻 Maple Tree，但可以通过 map 区间表理解范围映射的基本不变量。读源码时要问节点保存了哪些范围信息，查询如何找到覆盖区间，更新如何保持相邻关系，旧节点何时释放。

实验用例要包括插入到最前、插入到最后、端点相接、完全覆盖、跨多个区间、删除中间子段、删除不存在区间。报告要写清半开区间和闭区间的差别。半开区间下，前一个终点等于后一个起点不重叠；闭区间下，端点相等可能冲突。很多区间题的错误都来自端点语义没有先定下来。

\noindent\textbf{伙伴系统：总量够不代表连续块够。}

算法题常把内存分配当成免费，但系统里分配器本身就是算法。伙伴系统把连续页按 2 的幂次阶管理。分配某阶块时，若该阶没有空闲块，就向更高阶找，找到后逐层拆分；释放时，如果伙伴块同阶且空闲，就合并成高一阶块。这个机制强调连续性和对齐，而不是只看空闲页总量。

一个系统可能总空闲页很多，却没有足够大的连续块。比如很多小块分散在不同位置，总和超过请求大小，但无法合并成所需阶数。这就是碎片问题。刷题里可以把它类比为区间合并和二进制对齐，但不能简化成普通计数。每个空闲块必须只出现在它当前阶的链表里，拆分和合并都要维护这个不变量。

实验可以设总页数为 1024，维护从零阶到十阶的空闲链表。分配时打印从哪个高阶拆下，释放时打印伙伴编号和是否合并。构造一个碎片案例：先分配多个小块，再释放其中一部分，让总空闲量看似足够，但高阶请求失败。这个实验能把位运算、链表、区间、贪心分配和工程性能放到一起理解。

把这个视角带回 C++ 容器，读者会更容易理解为什么大量链表节点、哈希节点和树节点的分配成本不能忽略。复杂度都是线性或对数，连续数组可能因为缓存局部性和少分配而快很多；节点容器可能因为稳定身份和常数删除而必要。算法复杂度给出数量级，分配器和内存局部性决定真实表现。

\noindent\textbf{题解写作：每道题至少回答一个自己的问题。}

写题解时，不要让所有题都长成同一张表。可以有统一的复盘框架，但正文必须回答本题自己的问题。两数之和要回答为什么先查再插入；无重复子串要回答左端为什么不能回退；接雨水要回答何时能结算某个位置；柱状图要回答弹栈时算谁；LRU 要回答为什么节点里要保存 key；Dijkstra 要回答堆中旧状态为什么可以跳过；背包要回答容量循环方向从哪里来。

一个合格段落应当增加新信息。它可以解释一个错误实现为什么错，可以给一个最小反例，可以把暴力中的重复工作定位出来，可以说明一个容器的副作用，可以把源码结构和刷题模型做有限对照。如果一个段落只是在说“注意边界”“使用哈希优化”“复杂度降低”，但没有说明具体边界、具体 key、具体降低了哪一步，就应当重写。

题解还要有实验意识。每个高频题至少保留一个暴力对拍版本或小规模枚举器。前缀和题用随机数组对拍，二分题用线性枚举答案对拍，单调栈题打印入栈弹栈，DP 题打印小表，图题打印完整状态，链表题画节点连接。实验不是附属内容，它是把“我觉得对”变成“我知道为什么错不了”的过程。

最后是语言衔接。章节不是题号仓库，应该从一个问题过渡到下一个问题：连续内存引出双指针和窗口，窗口失效引出前缀和与单调队列，栈的最近未完成结构引出表达式和柱状图，动态候选集引出堆和最短路，子问题重复引出 DP，动态有序关系引出区间树和源码结构。这样的书读起来才像教材，而不是把许多答案贴在一起。

\noindent\textbf{源码与系统结构精读。}

刷题教材加入 Linux 源码，不是为了让读者背内核字段名，而是为了让读者看到同一类数据结构在真实系统中为何会变复杂。力扣题通常默认单线程、对象生命周期简单、内存分配随时可用；内核代码必须同时考虑并发读写、引用计数、缓存局部性、分配上下文和对象同时进入多个索引。下面这些小节仍然坚持单点分析，每节只追一个问题。

\noindent\textbf{list head：为什么节点嵌进对象。}

普通 C++ 初学者会把链表理解为“容器里存元素”。Linux 的侵入式链表不是这样。业务对象中嵌入链表节点，链表节点只负责前后指针，不拥有对象，也不知道对象类型。遍历时通过成员偏移从链表节点回到外层对象。这个设计看起来底层，但它直接解决三个工程问题：减少额外分配，让对象可以同时进入多条链表，保持对象地址稳定。

把它和 LRU 放在一起看会很清楚。刷题版 LRU 中，缓存条目既需要按 key 被哈希表找到，又需要按新旧顺序进入链表。如果链表节点和业务对象分离，删除和移动时要处理额外间接层；如果对象自己带链表节点，哈希表可以直接指向对象，链表只改变对象中的链接字段。内核里一个页面对象可能同时出现在 LRU 链、回写链或其他状态链上，所以对象中可以有多个链表节点字段，每个字段代表一条不同关系。

单点实验是写一个 Item，里面有 key、value、lru node、dirty node。先把三个 Item 挂入 LRU 链，再把其中两个挂入 dirty 链。删除一个对象的 dirty node，不应影响它在 LRU 链中的位置；释放对象前，必须先确认它已经从所有链和索引中摘除。这个实验比单纯反转链表更接近源码思维，因为它把“节点身份”和“对象生命周期”放在同一处观察。

\noindent\textbf{hlist：哈希桶为什么不一定用完整双链。}

哈希表有大量桶，很多桶可能为空。如果每个桶头都使用完整双向链表节点，桶数组本身就会消耗更多空间。hlist 的设计可以理解为桶头较轻，节点保存足够信息以支持已知节点删除。桶头只需要知道第一个节点，节点除了 next，还保存指向自己的前驱链接位置。这样删除已知节点时，可以直接修改前驱的 next 或桶头的 first，不必从桶头扫描找前驱。

刷题里哈希冲突通常被标准库隐藏，但很多题都在使用同一思想：key 先定位桶，桶内再比较实际 key。两数之和的 key 是历史值，字母异位词的 key 是规范签名，前缀和计数的 key 是前缀值，账户合并的 key 是邮箱。高质量题解要说清 key 和 value，不应只说“用 unordered map”。如果 key 设计错了，哈希表再快也会得到错误答案。

实验可以写两个固定桶数的小哈希表。第一版桶内普通单链表，删除已知节点必须找前驱；第二版节点保存前驱链接位置，删除时不用扫描。再设计一个遍历中删除的场景：如果读者正在走桶链，写者删除节点后什么时候能释放对象？单线程实验只需保存 next，内核并发场景则需要锁或 RCU。这个单点把哈希冲突、删除成本和并发生命周期自然连起来。

\noindent\textbf{rbtree：比较逻辑为什么交给调用者。}

C++ 的 map 把比较器和节点管理封装在容器里。Linux rbtree 更像一个底层平衡树工具：树节点嵌在业务对象里，插入和查找时，调用者自己沿树比较 key，找到位置后再调用库函数链接和修复平衡。这样设计的好处是业务对象可以按不同字段进入不同树，内存分配和对象生命周期也不被通用容器隐藏。

红黑树单点要抓住“旋转不改变中序顺序”。插入新节点后，树可能违反颜色和平衡约束，修复函数通过变色和旋转调整局部结构。旋转前后，中序遍历得到的 key 顺序相同，所以查找语义不变。很多教材一上来讲红黑树性质，读者很快迷路；源码阅读可以先问一个更小的问题：调用者如何比较 key，库函数如何在不破坏有序性的前提下修复平衡。

刷题对应的是动态有序集合。日程安排要找前驱后继，滑动窗口有序统计要插入删除并取极值，时间键值存储在每个 key 内部按时间二分，动态区间表需要按起点维护顺序。哈希表适合等值查找，不能回答前驱后继。实验用 map 实现日程安排，再手写一个简化 BST 插入，不要求实现红黑修复，只要求打印中序序列。等读者理解有序语义稳定后，再看平衡修复才不会把重点放错。

\noindent\textbf{XArray：稀疏整数索引不是普通哈希。}

XArray 面向整数索引到对象指针的映射，适合大而稀疏的索引空间。巨大数组会为空洞浪费空间，普通哈希表能做等值查找，但对按序推进、范围扫描、标记和并发读写支持不一定合适。XArray 这类结构把整数索引按位分层，只为实际存在的路径分配节点。它和 Trie 的直觉相似：字符串 Trie 每层消耗一个字符，整数分层索引每层消耗几位。

刷题里可以用三个模型建立直觉。位 Trie 用二进制位分层，常见于最大异或；坐标压缩把稀疏大值域映射到紧凑数组下标；有序 map 支持稀疏 key 的顺序遍历。XArray 和它们都不完全相同，但共享一个问题：key 空间很大，实际对象稀疏，需要在空间、查找、遍历和更新之间折中。

实验可以写一个简化二进制 Trie 保存整数集合，支持插入、查找、寻找不小于某个值的下一个元素。再和数组、哈希表、map 比较：数组适合小连续值域，哈希适合稀疏等值查找，map 适合动态有序，Trie 适合可按位分解的 key。读 XArray 时不要先背节点字段，先回答索引如何分层、空洞如何表示、读者什么时候能看到新对象。

\noindent\textbf{Maple Tree：范围映射的三个操作。}

范围映射比单点映射更容易出错。给定地址要找到覆盖它的区间，插入新区间要检查重叠，删除区间可能把原区间切开。Maple Tree 服务的正是这种范围索引场景。它比普通 map 复杂，是因为真实内核还要追求读路径性能、处理并发、减少分配和维护更复杂的节点布局。但教材阶段先抓三个操作：点查询、插入、删除。

用简化 map 区间表可以训练这些操作。点查询先找起点不大于地址的最后一个区间，再看终点是否覆盖；插入只需检查前驱和后继，因为已有区间不重叠且按起点有序；删除有四种情况：删除整段，截掉左侧，截掉右侧，从中间挖掉一段并分裂。每种情况都要画图，否则代码很容易在端点处错。

这和力扣 56、57、729、731 是一条线。合并区间是离线排序扫描，插入区间是一次性合并连续重叠段，我的日程安排是动态插入并检查冲突。Maple Tree 是系统级动态范围集合。读者不需要一开始理解全部源码，但要知道为什么一个普通哈希表无法回答“哪个区间覆盖这个地址”，也无法按地址顺序找下一个区域。

\noindent\textbf{RCU：删除要拆成摘除和释放。}

RCU 最适合从链表删除理解。单线程删除一个节点，可以改前后指针后立刻释放节点。读多写少并发场景不行：读者可能已经拿到旧指针，写者把节点从新索引中摘除后，旧读者仍可能继续访问它。于是删除必须拆成两个阶段：先从结构中摘除，让后来的读者找不到；再等待所有可能看到旧节点的读者离开；最后释放内存。

这个拆分和刷题里的懒删除有相似直觉，但目的不同。滑动窗口中位数的堆懒删除是因为 priority queue 不能删除任意元素；Dijkstra 的旧堆状态是因为没有 decrease-key；RCU 延迟释放是为了内存安全。相似点只是“逻辑不可用”和“物理存在”分离，不能把它们说成同一种机制。教材要防止概念类比过头。

实验可以模拟一个全局配置表。读者进入读侧临界区后读取当前指针；写者复制旧表，修改副本，发布新指针，把旧表放入延迟释放列表。只有当活跃读者计数为零时，才释放旧表。这个实验能让读者看到，读路径轻量的代价是写路径更复杂。回到算法题，读者也会更谨慎地使用“删除”这个词：从集合不可达、从答案中移除、从内存释放，三者不是同一件事。

\noindent\textbf{调度队列：为什么动态最小值不总是堆。}

调度器需要在可运行任务中选择下一个运行者。一个简化模型是每个任务有虚拟运行时间，每次选择虚拟运行时间最小的任务运行，运行后它的虚拟运行时间增加，再放回候选集合。直觉上堆可以取最小值，但真实调度还需要删除任意任务、处理睡眠唤醒、按顺序遍历、维护权重和调度类。堆只解决“取最小”，不解决全部操作。

刷题中类似问题很多。数据流中位数用双堆，因为只需要两半极值；日程安排用有序表，因为需要前驱后继；任务调度器可以用堆模拟，也可以用最高频公式；滑动窗口最大值用单调队列，因为候选只按窗口顺序过期。结构选择必须列出操作集合：插入、删除任意元素、取最值、找前驱后继、范围遍历、按时间过期。只看一个操作很容易选错结构。

实验写一个简化调度模拟器。第一版用 priority queue，每次任务运行后压入新状态，睡眠任务只能懒删除；第二版用 set，可以按迭代器删除任意任务，再重新插入更新后的任务。比较两版代码，读者会看到接口能力的差异。复杂度同为对数，维护语义却不同。这个单点能把堆、有序树和系统队列联系起来。

\noindent\textbf{页缓存和 LRU：真实淘汰不是一道缓存题。}

力扣 LRU 缓存的规则很干净：访问就移动到头部，容量满就淘汰尾部。页面回收里的 LRU 思想复杂得多。页面可能是文件页或匿名页，可能干净或脏，可能正在写回，可能被引用，可能属于不同内存区域。系统不能简单删除“最旧页面”，还要考虑 I/O 成本、回收收益和并发状态。

教材中引入这个差异，是为了防止读者把刷题结构绝对化。哈希表加双向链表是 LRU 的核心索引组合，但真实系统可能使用多条链表达状态分层，例如活跃和非活跃，文件和匿名。页面从一条链移动到另一条链，不仅是顺序变化，也是状态变化。删除页面也不是释放一个普通对象，还要处理引用、脏页写回和映射关系。

实验可以从简化两级 LRU 开始。把页面分成 active 和 inactive 两条链，访问 inactive 页面时提升到 active，周期性把 active 尾部降级，回收从 inactive 尾部开始。再加入 dirty 标记：脏页不能立即释放，需要先模拟写回。这个实验不追求复刻内核，只让读者看到“链表表达状态，哈希表达定位，标记表达附加约束”。再回到 LRU 题，读者会更理解为什么节点要稳定、为什么多个索引要同步。

\noindent\textbf{高频设计题精读。}

设计题比普通单题更容易露出模板化问题，因为它们通常不是一个数据结构就能完成，而是多个索引协同维护。每次读设计题，都要写出操作表：每个操作需要什么查询、什么更新、什么删除、什么顺序。下面几个题都按这个方式分析。

\noindent\textbf{LeetCode 460 LFU 缓存：频次和时间两个维度。}

LFU 缓存要求淘汰使用频次最低的项；如果频次相同，淘汰最久未使用的项。它比 LRU 多一个频次维度。暴力做法可以每次淘汰时扫描所有 key，找最低频次和最旧时间，复杂度太高。要做到常数时间，需要两个索引：key 到节点信息，frequency 到同频节点链表。还要维护当前最小频次 min freq。

节点信息至少包含 key、value、freq，以及它在对应频次链表中的位置。get 命中后，节点频次加一：先从旧频次链表删除，再加入新频次链表头部；若旧频次链表空且旧频次等于 min freq，min freq 加一。put 新 key 时，若容量满，从 min freq 对应链表尾部淘汰最旧节点。插入新节点的频次为一，因此 min freq 重置为一。

这题最容易错在只更新频次，不更新同频时间顺序；或者删除节点后忘记维护 min freq。实验序列要包含容量为零、容量为一、同频淘汰、get 后频次变化、put 已存在 key、低频链表清空。每步打印 min freq、每个频次链表顺序和 key 表大小。LFU 的核心不是“哈希加链表”四个字，而是频次索引和时间索引同时一致。

\noindent\textbf{LeetCode 895 最大频率栈：为什么 frequency 到 stack 足够。}

最大频率栈要求 pop 时返回出现频率最高的元素；若频率相同，返回最近压入的元素。暴力可以每次 pop 扫描所有元素，找最高频和最近时间。优化时观察两个维度：value 到当前频率，frequency 到达到该频率的元素栈。push 一个值时，它的频率从 f 变成 f+1，把这个值压入 group[f+1]，同时更新 max freq。

pop 时直接从 group[max freq] 弹出最近达到最高频的值，然后把它的当前频率减一。若 group[max freq] 变空，max freq 减一。为什么不需要从旧频率栈中删除这个值？因为一个值每次频率增加都会在新频率栈留下记录，这些记录对应历史时刻；pop 只会从当前最高频栈弹出合法的最近记录。频率减一后，低频栈中旧记录仍然代表它之前达到低频时的顺序，后续在那个频率层可能被使用。

实验序列要让同一个值多次 push，并和另一个值形成同频竞争。打印 value frequency、max freq、每个 group 栈。若实现使用单个堆，也能做，但要处理过期记录；frequency 到 stack 的结构更贴合题意，因为“同频最近”天然是栈。这个题训练的是把一个复杂排序规则拆成两级索引，而不是盲目使用优先队列。

\noindent\textbf{LeetCode 981 时间键值存储：按 key 分组再二分。}

时间键值存储支持 set(key,value,timestamp) 和 get(key,timestamp)，返回该 key 在不超过 timestamp 的最新值。暴力做法把所有记录放在一个列表里，每次 get 扫描。优化的第一步是按 key 分组，因为不同 key 的记录互不影响。对每个 key，保存按时间递增的 pair 列表。get 时在该 key 的列表中二分，找最后一个时间不超过目标的记录。

这题的关键条件是 set 的 timestamp 通常按递增顺序到来。如果不递增，就需要插入时保持有序，或者最后排序。二分查找的是 upper bound：第一个大于目标时间的位置，答案是它前一个。若位置为开头，说明没有可用记录。这里不能用哈希表直接按 timestamp 查，因为查询是“不超过”的前驱问题，不是等值问题。

实验要覆盖 key 不存在、时间早于第一条记录、等于某条记录、落在两条记录之间、晚于最后记录、多个 key 交错写入。报告里要写清为什么结构是 hash map 到 vector，而不是一个全局 map。这个题非常适合训练“先按独立维度分组，再在组内使用有序结构”的思想。

\noindent\textbf{LeetCode 632 最小区间：多路归并里的覆盖条件。}

最小区间覆盖 K 个列表，要求区间内至少包含每个列表一个数。暴力可以从所有元素中枚举左右端，再检查是否覆盖每个列表。优化从多路归并视角看：每个列表当前选一个元素，则当前 K 个元素天然覆盖所有列表，区间由其中最小值和最大值决定。为了缩小区间，应当推进当前最小值所在的列表，因为当前最大值不变时，只有提高最小值才可能缩短区间。

小顶堆保存每个列表当前元素，同时维护当前最大值。每次弹出最小元素，更新答案区间，然后把该元素所在列表的下一个元素压入堆，并更新最大值。如果某个列表耗尽，就无法再保持覆盖所有列表，算法结束。这个过程和合并 K 个有序链表相似，但多了当前最大值和覆盖区间的目标。

证明时看当前状态覆盖了所有列表。若不推进最小元素，而推进其他列表，当前最小值仍在，区间左端不变，最大值还可能变大，不可能得到更优区间。因此推进最小所在列表是安全的。实验覆盖某个列表只有一个元素、多个相同值、答案在开头、答案在中间、列表长度差异很大。C++ 中堆元素保存值、列表编号和下标，避免复制整个列表。

\noindent\textbf{LeetCode 815 公交路线：节点选错会爆炸。}

公交路线题要求从起点站到终点站最少乘坐几辆公交。若把站点作为节点，边表示两个站在同一路线上相邻，建图会非常大；若把同一路线内所有站两两连边，更会爆炸。更稳的建模是站点到路线列表的哈希表，BFS 的层次表示乘坐公交次数。起点站能乘坐的所有路线作为第一层路线，访问一条路线时，可以到达这条路线上的所有站，再从这些站扩展到其他路线。

这里的访问标记要分路线和站点两个层次。路线一旦乘坐过，就不需要再次乘坐；站点如果已经用于扩展过，也可以避免重复扫描它对应的路线列表。目标是到达终点站，不一定要构造完整的路线图。若起点等于终点，答案为零；若终点站不在任何可达路线中，返回失败。

实验要比较两种建模的扩展数量。构造多条长路线共享少量换乘站的输入，站点图会产生大量邻接关系，而路线 BFS 只在需要时通过站点找到可换乘路线。这个题说明图题第一步不是写 BFS，而是选择节点含义。节点选错，后面所有优化都只是在错误模型上补救。

\noindent\textbf{LeetCode 1235 规划兼职工作：排序、二分和 DP 合体。}

规划兼职工作给出若干工作，每个工作有开始时间、结束时间和收益，要求选择不重叠工作最大化收益。按收益贪心会失败，因为高收益工作可能阻塞多个组合后更高的工作。暴力枚举子集会指数爆炸。正确切入点是按结束时间排序，定义 dp[i] 表示考虑前 i 个按结束时间排序的工作时的最大收益。

对于第 i 个工作，有两种选择：不选它，收益是 dp[i-1]；选它，则要找到最后一个结束时间不超过当前开始时间的工作 j，收益是 dp[j] 加当前收益。因为工作已按结束时间排序，j 可以用二分找到。这个二分不是在答案空间上，而是在已经排序的结束时间数组上找前驱边界。

实验要覆盖完全不重叠、全部重叠、端点相接是否允许、同结束时间、同开始时间、收益很大的单个工作和多个小工作组合。打印排序后的工作、每个 i 找到的 j、选择和不选择的收益。这个题很好地展示了多技术衔接：排序制造可二分的结束时间，二分找到兼容前驱，DP 比较选或不选。任何一个环节单独背模板都不够。

\noindent\textbf{LeetCode 1675 最小偏移量：先规范化状态空间。}

数组最小偏移量的操作规则有点绕：偶数可以除以二，奇数可以乘以二。若直接搜索所有可能状态，分支很多。关键观察是每个奇数最多先乘二一次，变成偶数后就只能不断除以二；每个偶数也只能不断除以二。于是可以先把所有数规范化到“可往下减少的最大形态”：奇数乘二，偶数保持。之后每次只尝试降低当前最大偶数，同时维护当前最小值。

用大顶堆保存当前所有数中的最大值，另有变量保存当前最小值。每轮用最大值减最小值更新答案；如果最大值是偶数，就把它除以二放回堆，并更新最小值；如果最大值是奇数，就无法继续降低它，算法结束。为什么每次只动最大值？偏移量由最大和最小决定，降低非最大值不会降低当前最大值，反而可能降低最小值让偏移更大。

实验覆盖全奇数、全偶数、混合、已经最优、最大值很大、多个相同值。打印堆顶、当前最小值和答案。这个题的单点是“规范化状态空间”：先把所有数推到统一方向可操作的边界，再用堆做贪心推进。它和很多搜索题相通，先把可逆或双向操作改成单向可控过程，复杂度才降下来。

\noindent\textbf{实验报告样例库。}

下面给几份更具体的实验报告样例。它们不是为了增加形式，而是给读者一个可执行标准：每个实验都能复现一个算法为什么对，或者一个错误实现为什么错。

\noindent\textbf{报告样例：前缀和对拍。}

问题：验证和为 K 的子数组哈希版是否漏掉边界。暴力基线：两层枚举左右端点，内部用累加变量得到区间和。优化版本：扫描前缀和，哈希表保存历史前缀频次。输入生成：数组长度从零到二十，元素范围从负三到三，K 从负五到五随机选取。每组数据同时跑暴力和优化，结果不同就打印数组、K、当前前缀、哈希表内容和下标。

预期发现：若漏掉初始前缀零，所有从下标零开始的答案会少计；若先更新再查询，K 为零时可能多算空区间；若用滑动窗口替代哈希，含负数输入会失败。结论：本题依赖代数等式，不依赖单调性；哈希表 value 必须是频次，因为题目问数量；查询更新顺序是算法语义的一部分。

\noindent\textbf{报告样例：单调栈状态打印。}

问题：验证柱状图最大矩形宽度计算。暴力基线：每根柱子向左右找第一个更矮位置。优化版本：单调递增栈加末尾哨兵。输入集合：空数组、单柱、严格递增、严格递减、全相等、中间低谷。每次弹栈时打印当前下标、被弹出下标、高度、左边界、右边界、宽度和面积。

预期发现：若没有末尾哨兵，严格递增数组会漏算；若宽度写成 right-left，递减数组会面积过大；若把当前柱高度当作结算高度，中间低谷会错。结论：弹出时被弹出的柱子才是矩形高度，当前柱只提供右边界，弹出后的栈顶提供左边界。

\noindent\textbf{报告样例：二分谓词审查。}

问题：验证船运包裹容量二分。暴力基线：从最大包裹到总重量线性枚举容量，找到第一个可行容量。优化版本：左闭右开二分第一个可行容量。检查函数：按顺序装包，当前天放不下就开新天，统计天数是否不超过 D。输入集合：一天运完、每天一个包裹、最大包裹等于答案、多个相同重量、容量小于最大包裹的非法 mid。

预期发现：若下界不是最大包裹，检查函数必须额外处理单包超过容量；若天数统计初始值错，边界会偏一；若二分返回未验证的 left，也可能隐藏谓词错误。结论：答案二分的核心是可行性谓词单调，外层循环只是利用这个谓词找边界。

\noindent\textbf{报告样例：LRU 索引一致性。}

问题：验证 LRU 的哈希表和链表是否一致。暴力基线：用 vector 保存从新到旧的 key，每次操作线性移动。优化版本：哈希表加双向链表。操作序列：容量为二，依次 put 一、put 二、get 一、put 三、get 二、put 一更新、put 四。每步打印链表顺序、哈希表 key 集合、每个 key 指向的节点值。

预期发现：若 get 后不移动节点，淘汰顺序会错；若 put 已存在 key 时新建节点，链表会出现重复 key；若淘汰尾部时忘记删哈希表，后续 get 会访问失效节点。结论：LRU 的不变量有两个，哈希表到链表节点一一对应，链表顺序表达最近使用到最久未使用。

\noindent\textbf{报告样例：图状态维度。}

问题：验证障碍消除最短路是否错误合并状态。暴力基线：小网格上用 BFS 保存行、列、剩余消除次数。错误版本：visited 只保存行和列。输入集合：必须绕路保留消除次数的网格、起点终点相邻、K 为零、K 大于障碍数。每次入队时打印位置、剩余次数和距离。

预期发现：二维 visited 会把“到达同一格但剩余资源更多”的状态丢掉，导致错误不可达或路径过长。结论：BFS 第一次到达最短，只对完整状态成立。状态维度少了，访问标记就会错误剪枝。

\noindent\textbf{报告样例：背包循环方向。}

问题：验证 0-1 背包一维压缩方向。暴力基线：二维 dp，阶段是物品下标，容量是当前目标。错误版本：一维数组容量正序。输入集合：一个物品重量一、容量二；两个物品重量一和二、容量三；分割等和子集中的小数组。打印每轮物品处理后的 dp 数组。

预期发现：正序容量会让同一个物品在同一轮被重复使用。倒序容量读取的是上一轮阶段状态，符合每个物品只能用一次。结论：循环方向不是模板差异，而是题目对物品使用次数的约束。

\noindent\textbf{章节衔接重写原则。}

本书后续扩展时，章节应该按能力递进，而不是按题号堆叠。线性结构部分先讲连续内存、下标和缓存局部性，再讲双指针和窗口；窗口失效后，自然引出前缀和、哈希和单调队列。栈和队列部分先讲最近未完成结构，再讲单调栈、表达式解析和柱状图。堆部分先讲动态极值，再讲 Dijkstra、数据流中位数和多路归并。

树和图部分也要有顺序。树先讲遍历方向：自顶向下传约束、后序向上汇总、层序传播；再讲 BST 的有序性和 Trie 的前缀共享。图先讲节点和边的建模，再讲 BFS 的层次、DFS 的连通块、拓扑的依赖、并查集的集合代表，最后讲带权最短路。这样读者能看见一个概念如何扩展，而不是在相邻页里突然切换题型。

动态规划部分必须从暴力搜索树开始。线性 DP 解释最后一步，二维 DP 解释两个前缀，背包解释阶段和容量，区间 DP 解释最后合并点，状态机解释约束转移，状态压缩解释集合。每个小节都要有一个完整表格实验。空间压缩只能在完整状态讲清后出现，不能一开始就给滚动数组代码。

Linux 源码专题放在基础算法之后更合适。读者先掌握链表、哈希、有序树、区间、队列、堆、分配器这些算法结构，再去看内核里的侵入式节点、hlist、rbtree、XArray、Maple Tree、RCU、伙伴系统和 SLUB，才不会被字段名淹没。源码章节不是另一本操作系统书，而是把算法结构放进真实工程约束里重新审视。

目录也要克制。章标题保留大的学习路线，小节标题只在正文中服务阅读，不把每个微型点都塞进目录。一本教材的视觉舒适来自层级稳定、段落长度适中、代码块清晰、图示服务推理。章节过碎会让读者感觉像在翻题库；章节过长没有锚点又会迷路。后续排版应在这两者之间取平衡。

\noindent\textbf{新增专题候选清单。}

如果继续扩写，优先补这些不是模板、能提升质量的专题。第一，字符串算法专题：KMP 的失败函数为什么能跳转，Z 函数如何维护右边界，Manacher 如何把回文半径线性化，后缀数组和后缀自动机只作为进阶引导。每个算法都必须从暴力匹配出发，而不是直接给公式。

第二，数学和位运算专题：快速幂、欧几里得、质数筛、组合数、模逆元、同余、异或线性基。面试刷题不一定高频，但它们能训练代数不变量。写法上要强调适用条件，例如模逆元要求互质或模数为质数，位移要注意有符号整数。

第三，图高级专题：强连通分量、割点桥、最小生成树、二分图匹配、网络流只做求职向概览。重点不是铺算法名，而是说明它们分别回答什么问题：强连通回答有向互达，桥回答删边连通性，最小生成树回答连接全部点的最低总成本，匹配回答两侧资源配对。

第四，工程性能专题：缓存局部性、分配器、哈希退化、字符串拷贝、迭代器失效、递归栈、输入输出性能。每个点都要配小实验。算法书如果不讲这些，读者容易以为复杂度相同的代码真实表现也相同。

第五，面试表达专题：如何在五分钟内讲清一题，如何在追问中判断条件变化，如何承认不确定并回到暴力，如何用反例纠正错误贪心。这个专题能把刷题能力转成沟通能力，但必须用具体题演示，不能写成空泛建议。

\subsection{质量审查和错题复盘}

\noindent\textbf{反模板质量审查。}

这一节给后续写作和修订使用。算法教材最容易出现的低质量问题，不是少写几个题，而是把每道题都写成相同段落：先说暴力，再说优化，再说边界，再说 C++。这种结构本身可以作为复盘框架，但不能成为正文内容。正文必须回答本题自己的问题，否则读者读十页之后只会记住一套空话。

\noindent\textbf{审查一：每题是否有自己的核心疑问。}

每道题提交前，先写一句“本题真正要解释的问题是什么”。两数之和的问题是为什么历史表要先查再插入；无重复子串的问题是为什么左边界不能回退；接雨水的问题是何时能结算某一侧；柱状图的问题是弹栈时谁的左右边界确定；课程表的问题是边方向如何影响输出顺序；零钱兑换的问题是最少数量和组合数量的状态含义不同。若这句话写不出来，说明题解还停留在题型标签上。

核心疑问必须能被一个反例检验。比如左边界回退用 \code{abba} 检验，前缀和初始零用单元素等于目标检验，柱状图宽度用递减数组检验，背包循环方向用单物品容量二检验，图状态维度用剩余资源不同的同位置状态检验。没有反例的“注意边界”不是分析，只是提醒。

一个题可以有多个解法，但每个解法都要对应自己的疑问。接雨水的前后缀数组回答“每个位置的左右最高如何预处理”，双指针回答“何时能用较低已知边界结算”，单调栈回答“什么时候槽底知道左右边界”。如果把三种解法都写成“维护状态，线性扫描”，就是在抹掉差异。

\noindent\textbf{审查二：暴力版本是否真的服务推导。}

暴力版本不是形式。它应该说明候选空间是什么，检查条件是什么，复杂度慢在哪一步。三数之和暴力枚举三元组，慢在组合数量和重复答案；排序双指针正是利用单调性和相邻重复来减少这两类问题。编辑距离暴力尝试三种操作，慢在相同前缀对反复出现；DP 状态正是给前缀对命名。若暴力版本只写“暴力会超时”，它没有提供推导价值。

暴力也不一定要写代码，但必须能口头执行。读者应当能拿一个小输入，按暴力过程跑完，看到哪些操作被重复。窗口题可以枚举所有左右端；图题可以从每个起点重新搜索；区间题可以两两比较；设计题可以每次操作扫描全部元素。只有低效过程清楚，优化才不是突然出现。

审稿时可以问一句：优化到底替换了暴力中的哪一行工作？哈希替换历史扫描，前缀和替换重复区间求和，单调栈替换左右边界寻找，堆替换每轮找极值，DP 替换重复递归，并查集替换重复连通搜索。若回答不上来，就说明题解在堆算法名。

\noindent\textbf{审查三：状态是否有不变量。}

高质量题解必须写出状态含义，而不是只写变量名。slow 左侧保存什么，窗口内满足什么，哈希表保存当前下标之前还是包含当前下标，栈里是等待右边界的柱子还是等待更高温的日子，堆顶是否可能过期，dp[i][j] 表示哪个前缀，parent 数组是不是原图路径。这些问题都应在正文中说清。

不变量要能指导代码顺序。前缀和题先查再更新，是因为哈希表表示历史前缀；颜色分类遇到二不移动扫描指针，是因为换回来的元素仍属于未知区；LRU 更新已有 key 要移动节点，是因为访问也改变最近使用顺序；拓扑排序入度变零才入队，是因为所有先修已满足。若状态含义不能决定代码顺序，说明它还不够具体。

不变量还要能解释结束条件。二分结束时 left 为什么是答案，BFS 第一次到达为什么最短，单调队列队首为什么是当前窗口最大值，DP 表最后一个格子为什么是全局答案，回溯到叶子时为什么可以收集路径。很多错误代码不是主体循环错，而是答案读取时机错。

\noindent\textbf{审查四：C++ 内容是否和题目相关。}

C++ 提醒不能每题都写同一句“注意 \code{std::int64_t}”。只有涉及求和、乘积、面积、距离、时间和路径代价时，才重点讨论溢出。只有哈希查询语义会影响状态时，才强调 \code{find} 和 \code{operator[]} 的区别。只有保存节点身份时，才强调 \code{list} 迭代器稳定。只有比较器决定排序不变量时，才重点检查严格弱序。

容器选择也要跟访问模式绑定。vector 适合连续扫描和随机访问，deque 适合两端操作，unordered map 适合等值历史查询，map 适合前驱后继，priority queue 适合动态极值，list 适合已知节点常数移动，array 适合固定字符集计数。若题解只说“用某容器更方便”，没有说明它回答哪个查询，就应该重写。

代码块要少而精。一本教材不需要每道题都贴完整最终代码；更重要的是给关键骨架、状态更新和易错行。代码排版必须稳定：缩进清晰，命名表达状态，辅助函数拆出检查逻辑，复杂条件不要塞在一行。代码服务推理，不是用来填满页面。

\noindent\textbf{审查五：Linux 源码对照是否克制。}

源码对照只能在确实能解释访问模式时出现。LRU 对照 list head，动态区间对照 Maple Tree，哈希桶对照 hlist，动态有序集合对照 rbtree，读多写少删除对照 RCU，分配连续页对照伙伴系统。这些对照有价值，因为它们回答“真实系统为什么比题目复杂”。如果只是写“Linux 里也用了某结构”，没有说明对象生命周期、并发、分配或多索引，就没有教学价值。

源码字段名会随版本变化，教材不应把字段名当重点。重点应是接口语义、对象关系、热路径和不变量。读 list head 就问节点如何嵌入对象；读 rbtree 就问比较逻辑在哪里；读 XArray 就问稀疏整数索引如何分层；读 Maple Tree 就问范围查询如何保持相邻关系；读 RCU 就问删除为何不能立刻释放。每次只追一个问题，读者才不会迷路。

源码实验也要可执行。写最小侵入式链表，写固定桶数哈希表，写 map 区间表，写简化伙伴分配器，写读写配置表模拟 RCU。实验不复刻内核，而是复现一个不变量。报告里必须写明和刷题题目的相同点和差异点，防止类比过度。

\noindent\textbf{审查六：重复句子必须删除。}

修订时可以直接搜索一些危险句式：这题用某结构、注意边界、时间复杂度降低、维护状态、不要套模板、实验时对拍、C++ 要注意容器。这些句子不是绝对不能出现，但如果没有跟具体题目绑定，就应删除或改写。比如“注意边界”要改成“目标串含重复字符时 formed 表示达标字符种类，不是匹配总数”；“维护状态”要改成“哈希表保存当前下标之前每个前缀和的出现次数”。

同一章节内，相邻题目的开头也要避免同构。可以用题面冲突开头，可以用反例开头，可以用暴力慢点开头，可以用错误实现开头，可以用源码类比开头。读者需要节奏变化。若十道题都以“先从暴力开始”开头，即使内容不同，也会产生复制感。

真正必要的重复是术语稳定，而不是段落重复。比如“不变量”“候选”“历史状态”“前驱后继”“过期状态”这些词可以反复出现，因为它们构成学习框架；但每次出现都要落到不同对象上。稳定词汇让书有体系，重复段落会让书像拼接稿。

\noindent\textbf{审查七：章节是否能串起来。}

每章开头要告诉读者上一章解决了什么问题，本章接着解决什么问题。数组讲连续内存和下标，窗口讲连续区间的增量维护，前缀和讲窗口失效后如何用代数历史状态补救，单调结构讲如何保存未解决候选，堆讲动态极值，图讲状态扩展，DP 讲重复子问题，源码讲真实工程约束。这条线如果断了，读者就会觉得每章都在换话题。

章节结尾要给迁移任务，而不是只总结。学完前缀和，让读者比较 560、523、525、862；学完单调栈，让读者比较每日温度、柱状图、下一个更大元素；学完背包，让读者比较 416、494、518、322；学完图，让读者比较岛屿、课程表、网络延迟、公交路线。迁移任务能强迫读者看到条件变化，而不是只记题号。

目录层级也要服务这种串联。大章标题表达学习阶段，小节标题表达核心问题，过碎的小节不要全部进入目录。正文中的图示、表格和实验框可以帮助定位，但不应把每个提醒都变成标题。一本书的排版舒适，来自稳定层级和连续叙述，而不是标题数量。

\noindent\textbf{审查八：如何判断一段应该扩写还是删除。}

一段文字如果回答了为什么、如何证明、哪里会错、如何实验、如何迁移、如何和源码差异化，就值得扩写。比如“Dijkstra 堆里有旧状态”值得扩，因为它解释了 priority queue 的接口限制和正确性检查；“窗口要维护计数”只有在说明 formed、need、重复字符时才值得扩；“DP 可以空间压缩”只有在说明旧值来源时才值得扩。

一段文字如果只是换词重复前文，就应该删除。比如已经在题卡里讲过“哈希保存历史”，后面再写“哈希能快速查询历史”就没有新信息，除非加入具体 key、value、查询顺序或反例。已经在章节里讲过“数组连续内存缓存友好”，后面再重复这句话也没有意义，除非加入 vector 扩容、迭代器失效或和 list 的实验对比。

扩写和删除都要服务读者。字数目标是最低体量，不是写作目标。一本三十万字的教材如果每一万字都能教会一个新判断，就是扎实；如果只是同一句话重复三十遍，就是负担。后续修订时宁愿删除一千字模板，也要补一千字能解释一个具体错误的内容。

\noindent\textbf{审查九：最终交付前的抽样方法。}

最终交付前，不要只看总字数和编译成功。随机抽十道题，遮住标题之外的内容，只读正文第一段，判断是否能看出本题自己的问题。若第一段换到另一道题也能成立，就说明太模板。再抽五个代码块，检查缩进、字体、行宽、注释和变量名。再抽五个实验，确认它们能复现一个错误或证明一个不变量。

还要抽查跨章节引用。接雨水在数组窗口章节讲过，在单点深挖里又讲，两个地方应该互补，而不是重复。LRU 在链表、设计题和 Linux 对照里都出现，三个地方应分别强调链表移动、索引一致性和侵入式节点。Dijkstra 在图章节和设计实验里出现，应分别强调非负边证明和旧堆状态。重复题可以出现，但角度必须不同。

最后检查手机阅读体验。段落不能太碎，代码不能横向溢出，目录不要塞满细碎小标题，彩色强调不能影响正文黑色阅读，图示要解释状态流动。手机上读算法书最怕目录混乱和代码挤压，所以构建 PDF 和 EPUB 后都要实际打开检查。技术内容和排版体验必须一起达标。

抽查结果要写进修订记录。发现一个模板段，就回到对应题目重新补模型、反例和实验，而不是只换几个词。

\noindent\textbf{审查十：术语统一但例子必须变化。}

全书可以统一使用一些核心术语，例如候选、状态、不变量、历史信息、支配关系、边界、前驱、后继、过期状态、逻辑删除、物理释放。这些词能让读者在不同章节之间建立联系。但术语统一不等于例子重复。每次出现“不变量”，都要落到具体对象：窗口题是不重复窗口，数组题是已写入前缀，栈题是仍未结算的下标，DP 题是已经计算完的依赖状态，源码题是对象已经从索引摘除但还没释放。

同样，每次出现“支配关系”也要换成题目自己的语言。单调队列里，新下标前缀和更小且位置更靠后，所以支配旧下标；滑动窗口最大值里，新值更大且更晚，所以支配旧值；柱状图里，当前更矮柱不是支配旧柱，而是触发旧柱结算；贪心题里的支配常常要靠交换证明。若这些差异没有写出来，术语就会变成空壳。

最后抽查一句话是否合格，可以问它能否直接指导代码。比如“维护历史状态”不合格；“哈希表保存当前下标之前每个前缀和出现次数，扫描当前前缀时先查询再更新”合格。比如“注意链表指针”不合格；“改 current next 前先保存 next，反转后 previous 指向已处理前缀头，current 指向未处理第一节点”合格。教材中的每个关键句都应该尽量达到这种可执行程度。

\noindent\textbf{最终自测题。}

自测一：给一个含负数数组，问和至少为 K 的最短子数组。请说明为什么普通滑动窗口失效，为什么前缀和单调队列成立，队尾弹出条件是什么，队首弹出条件是什么。

自测二：给一个动态日程表，要求插入区间并拒绝冲突。请说明为什么哈希表不合适，为什么只检查前驱和后继，半开区间下端点相等是否冲突。

自测三：给一个带正权图，要求从起点到所有点的最短时间。请说明为什么 BFS 不够，Dijkstra 的过期堆状态如何处理，距离数组使用什么类型。

自测四：给一个链表设计题，要求 O(1) 移动节点。请说明节点由谁拥有，哈希表保存什么，删除节点时哪些索引要同步更新，是否能用 vector 保存节点。

自测五：给一个背包题，要求统计方案数。请说明物品能用几次，方案是否区分顺序，容量循环方向是什么，\code{dp[0]} 为什么初始化为一。

这些自测题不要求背出固定代码，但要求能完整说出模型、暴力、瓶颈、状态、不变量、边界和 C++ 风险。如果说不完整，就回到对应章节重写复盘报告。

\noindent\textbf{错题复盘样例。}

\noindent\textbf{错题一：窗口左端回退。}

错误现象：无重复字符最长子串在输入 \code{abba} 上返回三。错误原因通常是看到重复字符后直接把左端设为上次位置加一，没有检查上次位置是否仍在当前窗口内。处理最后一个 \code{a} 时，它上次出现在下标零，但当前窗口左端已经在下标二；如果把左端拉回一，就会让已经排除的字符重新进入窗口，破坏窗口无重复不变量。

复盘写法：窗口左端只能右移，不能回退。更新公式应当是取当前左端和上次位置后一位的较大值。这个错误说明我们没有把“窗口保存当前无重复子串”落实成不变量，只是在机械处理重复字符。修复后，再用 \code{abba}、\code{tmmzuxt}、全相同字符和全不同字符测试。报告中要记录每一步左端变化，而不是只写最终答案。

\noindent\textbf{错题二：前缀和漏掉开头区间。}

错误现象：和为 K 的子数组在输入 \code{[3]}、目标 \code{3} 时返回零。错误原因是哈希表没有预先放入前缀和零的一次出现。区间从下标零开始时，它对应的历史前缀是空前缀；如果不记录空前缀，就等于告诉算法“从开头开始的区间不存在”。

复盘写法：前缀和题必须解释历史状态的时间点。扫描到当前前缀 \code{s} 时，答案增加历史前缀 \code{s-k} 的出现次数；空前缀是合法历史状态，位置在数组开始之前。先查询再更新当前前缀，是为了避免把当前前缀当成历史前缀。修复后测试目标为零、负数、重复前缀和从开头开始的多段答案。

\noindent\textbf{错题三：二分检查函数错误。}

错误现象：船运包裹二分循环看似标准，但某些用例返回容量小于最大包裹。错误原因不是二分边界，而是检查函数没有在当前包裹超过容量时立即判不可行，或者下界没有设为最大包裹。二分只能在正确谓词上找边界，不能修复错误谓词。

复盘写法：先写谓词定义，再写二分。容量 \code{cap} 可行，含义是按原顺序把包裹分成不超过 D 段，每段和不超过 \code{cap}。如果某个包裹重量已经超过 \code{cap}，谓词必假。下界设为最大包裹，上界设为总重量，可以让检查函数更简单。每次失败用例都要打印 \code{cap} 和需要天数，确认单调性是否存在。

\noindent\textbf{错题四：柱状图宽度少一。}

错误现象：柱状图最大矩形在递减数组或含哨兵用例上面积偏小。错误原因常在弹栈后的宽度计算。弹出柱子后，当前下标是右侧第一个更矮位置；弹出后新的栈顶是左侧第一个更矮位置。有效宽度不包含左右两个更矮边界，因此是 \code{right - left - 1}。如果栈为空，说明左侧没有更矮边界，宽度是当前下标。

复盘写法：单调栈弹出时结算的是被弹出的柱子，不是当前柱子。当前柱子只是右边界。每一次弹出都要能说出左边界、右边界、高度和宽度。修复后测试严格递增、严格递减、全相等、单个柱子和末尾哨兵。这个错误说明我们背了栈模板，却没有理解弹出时机。

\noindent\textbf{错题五：背包循环方向写反。}

错误现象：分割等和子集在只有一个物品时却能凑出两倍容量。原因是 0-1 背包一维压缩时容量正序遍历，导致当前物品刚更新出的状态又被当前物品再次使用。这个错误在小样例上可能不明显，但一旦构造单物品反例就会暴露。

复盘写法：循环方向来自物品能使用几次。0-1 背包每个物品只能用一次，所以容量倒序，读取上一轮物品阶段的状态；完全背包物品可以无限使用，所以容量正序，允许当前物品在同一轮继续贡献。报告里不要写“背包模板”，要写“当前题每个数只能选一次，因此倒序容量”。再用零钱兑换 II 对比，说明组合计数为什么外层遍历硬币。

\noindent\textbf{错题六：图状态维度缺失。}

错误现象：带钥匙、障碍消除或剩余资源的网格最短路返回过大或错误不可达。原因是访问标记只按位置记录，把“同一位置但资源不同”的状态合并了。位置相同不代表未来能力相同；剩余障碍次数更多或钥匙集合不同，会改变后续可走边。

复盘写法：图题先定义状态，不只是节点位置。若未来决策依赖资源，就把资源放进状态和访问标记。障碍消除题的 visited 至少应包含行、列、剩余消除次数；钥匙迷宫要包含行、列、钥匙集合。BFS 的第一次到达最短，只对完整状态成立，不对被错误压缩的位置成立。测试要构造必须绕路保留资源的用例。

\noindent\textbf{错题七：LRU 更新已有 key。}

错误现象：LRU 在 \code{put} 已存在 key 后淘汰错误元素。原因是只更新了 value，没有把节点移动到最新位置，或者新建了一个重复节点而没有删除旧节点。LRU 的“使用”包括 get，也包括更新已有 key；只要访问或更新成功，它就应成为最近使用。

复盘写法：LRU 有两个不变量：哈希表中每个 key 对应链表中唯一节点，链表顺序从最新到最旧。更新已有 key 时，不能增加容量，不能产生重复节点，必须移动到头部。容量满淘汰时，要从链表尾找到 key，再从哈希表删除。测试容量为一、重复 put、get 后再 put、更新后淘汰等序列。

\noindent\textbf{错题八：Trie 搜索没有恢复访问标记。}

错误现象：单词搜索在某些路径上漏答案，或者同一格子被重复使用。原因通常是 DFS 进入格子后没有在返回时恢复访问标记，导致其他分支误以为该格不可用；或者根本没有访问标记，导致一条路径中重复使用同一格。Trie 剪枝只能判断前缀是否存在，不能替代网格路径的访问约束。

复盘写法：回溯代码必须成对出现：做选择、递归、撤销选择。访问标记属于路径状态，而不是全局永久状态。找到一个单词后可以清空 Trie 节点上的单词标记避免重复答案，但不能清空网格访问规则。测试棋盘有重复字母、同一单词多条路径、单字符单词和路径必须回退的用例。

\noindent\textbf{最终质量检查表。}

交付一本算法教材前，逐章检查以下问题。第一，是否每个核心算法都有暴力基线；第二，是否解释了优化替换了哪一步；第三，是否给出至少一个反例或边界实验；第四，是否说明 C++ 容器选择和风险；第五，是否有 Linux 或工程源码视角帮助读者迁移；第六，是否避免同一句话在大量题卡中机械重复。

对单题检查也用同样标准。题面模型是否一句话说清；输入规模是否决定复杂度预算；状态是否足够；不变量是否能反驳错误实现；边界是否进入测试表；代码是否有溢出、迭代器失效、默认插入、递归深度或比较器风险；变体是否重新证明。只有这些都通过，才算高质量题解。

最后检查阅读体验。章节标题不要过碎，小节应服务于连续讲解；目录只保留大结构，不把每个细碎点都塞进去；代码块使用稳定宽度和清晰缩进；正文以黑色为主，颜色用于提示、强调和代码；图示用于解释状态流动，而不是装饰。一本教材的质量不仅取决于内容量，也取决于读者能否顺着推理读下去。

\noindent\textbf{读者自查问答。}

\noindent\textbf{看到新题先问什么。}

第一问输入规模。规模决定暴力能否作为最终答案，也决定是否需要对数、线性、平方或伪多项式方法。第二问数据条件。数组是否有序，元素是否全正，字符集是否固定，图边权是否相同，树是否是搜索树，区间端点是闭区间还是半开区间。第三问输出目标。是判断、计数、最值、路径、所有方案还是修改后的结构。输出目标不同，保存的状态就不同。

第四问能否从暴力中看出重复工作。重复扫描历史值，可能用哈希；重复求区间和，可能用前缀；重复找极值，可能用堆或单调队列；重复计算子问题，可能用 DP；重复判断连通，可能用并查集。第五问优化依赖的条件是否稳定。如果条件变了，比如正数变成有负数，有序变成无序，等权变成带权，原算法就要重新证明。

\noindent\textbf{如何判断自己不是背模板。}

把答案遮住，只保留题面，尝试写出暴力版。如果暴力版说不清，说明还没理解候选空间。再尝试说出瓶颈是哪一步，而不是直接说某个算法名。然后说明优化结构维护什么状态，为什么这个状态足够回答输出目标。最后拿一个反例解释错误写法为什么错。能完成这四步，才说明不是背模板。

另一个检查方法是改条件。把数组加入负数，把返回一个答案改成返回所有答案，把判断存在改成统计数量，把静态输入改成动态更新，把普通树改成 BST 或把 BST 改成普通树。若你能判断原解法哪些部分保留、哪些部分失效，就说明理解的是结构。若只能说“这个题我见过”，说明还停留在题号记忆。

\noindent\textbf{如何把代码写得更像工程答案。}

工程答案要把风险前置。函数入口检查不可行条件，变量名表达状态语义，复杂判断拆成辅助函数，容器选择和访问模式一致。二分答案把检查函数单独写；图算法把邻居生成集中；链表操作把摘除和插入封装；DP 把初始化和转移分开。这样写不是为了增加行数，而是让不变量有清晰位置。

还要避免隐式成本。不要在热循环里反复构造大字符串；不要在 range-for 中复制复杂对象；不要把可能溢出的乘法留在 int；不要保存会被 vector 扩容失效的指针；不要用 \code{operator[]} 做纯查询；不要写不满足严格弱序的比较器。面试时间有限，也应主动说明这些风险，哪怕代码里只写最核心版本。

\noindent\textbf{如何读 Linux 源码不迷路。}

读源码前先写问题，不要漫无目的打开文件。比如“链表删除怎么保持前后指针”“红黑树插入时比较逻辑在哪里”“XArray 如何把整数索引分层”“RCU 为什么延迟释放”。每次只追一条路径，画出对象、索引、插入、查找、删除和释放。字段名可以暂时不背，但对象关系必须画清楚。

读完以后找一个刷题类比。list 对照 LRU，rbtree 对照动态区间，XArray 对照稀疏索引，RCU 对照逻辑删除和物理删除分离，伙伴系统对照分层空闲块。类比不是说二者完全相同，而是帮助理解访问模式。真正的源码能力来自差异：内核多了并发、分配上下文、引用计数和生命周期。

\noindent\textbf{如何长期复习。}

每周固定回看错题，而不是只做新题。错题报告按八项写：题面模型、暴力基线、瓶颈、优化结构、正确性、边界、C++ 风险、变体。一个月后随机抽十道题，只看题名重讲推导。如果讲不出暴力和瓶颈，就回到原章节；如果讲不出边界，就补实验；如果讲不出 C++ 风险，就补容器实验。

复习时把题按能力分组。前缀和组、单调结构组、答案二分组、拓扑和最短路组、背包组、区间 DP 组、设计题索引组合组。每组挑一题作为代表，再列两道变体。这样复习能建立迁移能力，而不是只增加题号数量。最终目标是看到新题时能自己搭桥：从题面条件到暴力，从暴力到瓶颈，从瓶颈到结构，从结构到证明。
